% 本文件由 LaTeX 排版 Agent 根据有效 translation.json 自动生成。
% author=Lactantius; work_id=0702; title=神圣训导撮要

\chapter{\WorkTitle{神圣训导撮要}}

% structure: p0001 source=h1 target=omitted reason=page-title-already-rendered-by-parent

致其兄弟彭塔狄乌斯\par

% structure: p0003 source=h2 target=section reason=relative-heading-rank; base=h2

\section{序言——全篇撮要及整部训导的计划与宗旨}

尽管我们很久以前为阐明真理与宗教而撰写的《神圣训导》诸卷，可能如此预备和塑造读者的心灵，以致其篇幅不致引起厌恶，其丰富不致成为负担；然而，你，兄弟彭塔狄乌斯，希望我为你撰写一部撮要，我想其原因是：我可以写些东西给你，并且你的名字可以因我的作品（姑且如此称之）而闻名。我将满足你的愿望，虽然将七卷巨著所论述的内容压缩在一卷书之内似乎是一项困难的事。因为当如此大量的题材被压缩在狭小的空间时，整个内容变得不够充实；而且因其过于简短而变得不够清晰，特别是许多为阐明论据所必需的论据和例证必然被省略，因为它们数量如此之多，以至单凭它们本身便足以构成一卷书。当这些被删除后，还有什么显得有用、什么显得明了呢？但我将尽可能在主题允许的范围内，既压缩散漫之处，又缩短冗长之篇；然而在这部旨在揭示真理的著作中，既不能因丰富而缺乏素材，也不能因理解而缺乏清晰。\par

% structure: p0005 source=h2 target=section reason=relative-heading-rank; base=h2

\section{第一章 论天主眷顾}

首先产生一个问题：是否存在天主眷顾，祂创造了或治理着这个世界？对此无人怀疑，因为几乎所有哲学家（除伊壁鸠鲁学派外）都众口一词地认为，世界不可能没有设计者而被造，也不可能没有统治者而存在。因此，伊壁鸠鲁不仅被最有学问的人驳倒，而且被所有凡人的见证和感知所驳倒。因为当人看到天地如此布置，万物如此安排，以至于它们不仅极其适宜于奇妙的美丽和装饰，而且也适宜于人的使用和其他生物的便利时，谁能怀疑天主眷顾的存在呢？因此，那按计划存在的事物，不可能没有计划而开始：这样，天主眷顾的存在便是确定的。\par

% structure: p0007 source=h2 target=section reason=relative-heading-rank; base=h2

\section{第二章 论天主独一，不能有多个}

接着有一个问题：天主是一位还是多位？这个问题确实包含很大的歧义。因为不仅个人之间意见不同，而且民族和邦国之间也不同。但若有人遵循理性引导，就会明白：主只能有一位，圣父也只能有一位。因为如果创造万有的天主也是主和圣父，那么祂必须是独一的，使同一位成为万物的元首和根源。除非万物都归于一位，除非有一位掌舵，除非有一位执缰，并且仿佛只有一个心智指导身体的所有肢体，否则世界便不可能存在。如果在蜜蜂的蜂群中有许多蜂王，它们便会灭亡或四散，而\par

\Quote{纷争以巨大的骚动攻击众王。}\par

如果在畜群中有几个头领，它们会争斗直到一个取得统治。如果在军队中有许多指挥官，士兵们便无法服从，因为发出了不同的命令；他们自己也难以维持统一，因为各人凭自己的心意谋取私利。因此，在这个世界的共同体中，除非有一位统治者，同时也是创造者，要么这个整体会瓦解，要么它根本不能被组合起来。\par

此外，全部权威不可能存在于多位神祇中，因为他们各自维持自己的职责和特权。因此，他们中没有一位能被称为\Quote{全能者}——这是天主的真正称号，因为他只能完成依赖于自身的事，而不敢尝试依赖于他人的事。武尔坎不会为自己要求水，尼普顿不会要求火；刻瑞斯不会要求技艺，密涅瓦不会要求果实；墨丘利不会要求兵器，玛尔斯不会要求竖琴；朱庇特不会要求医术，埃斯科拉庇俄斯不会要求雷霆：他宁愿忍受别人投来的雷霆，也不愿自己挥舞。因此，如果个体不能做万事，他们便拥有较少的力量和较少的权能；但那能完成整体、而非只能完成整体中最微小部分者，才应被视为天主。\par

% structure: p0012 source=h2 target=section reason=relative-heading-rank; base=h2

\section{第三章 诗人们关于独一天主的见证}

因此，只有一位天主，完美的，永恒的，不朽的，不能受苦的，不受任何环境或权能支配的，祂自己拥有一切，统治一切，人的心智既不能以思想估量祂，口舌也不能以言语描述祂。因为祂过于崇高伟大，无法被人的思想或言语所领悟。简言之，且不提先知——独一天主的宣讲者，就连诗人们、哲学家们和有灵感的妇女们也开口为天主的合一作证。俄耳甫斯提到那超越的天主，祂创造了天和太阳以及其它天体，创造了地和海洋。此外，我们自己的马罗时而称至高天主为一种神灵，时而称为一种心智，并说它犹如渗入肢体，促使整个世界的身体运动；又说天主渗透了天的高度、海与地的区域，一切生灵都从祂获得生命。就连奥维德也知道世界是由天主预备的，他时而称天主为万物的建造者，时而为世界的制造者。\par

% structure: p0014 source=h2 target=section reason=relative-heading-rank; base=h2

\section{第四章 哲学家们关于天主合一的见证}

但让我们转向哲学家，他们的权威被认为比诗人更为可靠。柏拉图断言祂的独一统治，说只有一位天主，世界由祂以奇妙的秩序预备并完成。他的门徒亚里士多德承认有一位理智统治着世界。安提司泰尼说，有一位神是本性上的神，是整个宇宙体系的治理者。要逐一列举泰勒斯、在他之前的毕达哥拉斯和阿那克西美尼，或者后来斯多葛派的克莱安特、克吕西波和芝诺，以及我们同胞中追随斯多葛派的塞内卡和图利乌斯本人关于至高天主的陈述，将是一项冗长的工作，因为所有这些人都试图界定天主的本质，并肯定世界由祂独自统治，且祂不从属于任何本性，因为一切本性都源于祂。\par

赫耳墨斯因他的德行和众多技艺的知识，配得\OriginalTerm{三倍伟大的}{Trismegistus}之名；他在学说的古老性上先于哲学家，并被埃及人尊为神。他在颂扬独一天主的威严时，用无尽的赞美称祂为主和父，并说祂没有名字，因为祂不需要名字，既然祂是独一的，而且祂没有父母，因为祂由自己并借自己而存在。他在给他儿子的信中这样开头：\Quote{理解天主是困难的，用言语描述祂是不可能的，即使对那些可能理解祂的人也是如此；因为完美者不能被不完美者理解，无形者不能被有形者理解。}\par

% structure: p0017 source=h2 target=section reason=relative-heading-rank; base=h2

\section{第五章——女先知们，即西比拉，宣告只有一位天主。}

现在要谈女先知们。瓦罗说曾有十位西比拉：第一位是波斯人，第二位是利比亚人，第三位是德尔斐人，第四位是基默里人，第五位是厄立特里亚人，第六位是萨摩斯人，第七位是库迈人，第八位是赫勒斯滂人，第九位是弗里吉亚人，第十位是提布尔提纳人，又名阿尔布内亚。他说其中只有库迈人的三卷书包含罗马的命运，被视为神圣；但几乎所有其他西比拉的书卷也存在，且通常被视为独立的，但它们被称作西比拉书，如同一个名字，只是那位据说生活在特洛伊战争时代的厄立特里亚人，在书中署上了自己的名字：其他人的著作则混在一起。\par

我所提到的这些西比拉，除了库迈人（她的书只有十五人委员会有权阅读）之外，都见证只有一位天主，祂是统治者、创造者、父亲，非由任何者所生，而是源于自身，祂从永远存在，并要到永远；因此惟有祂配得一切有生命之物的敬拜、敬畏和尊崇。我略去了她们的见证，因为我无法缩短它们；但若你希望看到，必须自己查阅那些书卷。现在我们继续余下的主题。\par

% structure: p0020 source=h2 target=section reason=relative-heading-rank; base=h2

\section{第六章——既然天主是永恒不朽的，祂不需要性别和继承。}

因此，如此众多而伟大的见证清楚地教导：世界中只有一个政府，一个权能，其根源无法想象，其力量无法描述。因此，那些想象诸神是从婚姻中生出的人是愚蠢的，因为性别本身及其结合，乃是天主赐予有死之人，为使每一族类通过后裔的继承得以保存。但不朽者何需性别或继承呢？因为既不涉及享受，也不涉及死亡？因此，那些被算作神明的，既然显然他们作为人出生，且生养了他人，他们分明是有死的。但他们被相信是神，因为当他们作为伟大而有力的君王时，因他们赐予人的恩惠，死后当得神性的尊荣；为他们建造庙宇和雕像，他们的记忆被保留并作为不朽者被颂扬。\par

% structure: p0022 source=h2 target=section reason=relative-heading-rank; base=h2

\section{第七章——论赫拉克勒斯的邪恶生活与死亡。}

虽然几乎万民都确信他们是神，但他们的行为，正如诗人和历史学家所记载，表明他们是人。谁不知道赫拉克勒斯生活的时代呢？因为他曾与阿尔戈英雄一同航海远征，又攻陷特洛伊，因普里阿摩的父亲拉俄墨冬的背信而杀了他。从那时算起，大约一千五百多年。据说他甚至不是由合法婚姻所生，而是由阿尔克墨涅通奸所出，且他自己也沾染了父亲的恶习。他从未远离女人或男色，走遍全世界，与其说是为了荣耀，不如说是为了情欲，与其说是为了杀戮野兽，不如说是为了生育儿女。虽然他战无不胜，却唯独被翁法勒征服，向她交出了他的棍棒和狮皮；穿着女装，匍匐在女子脚下，接受她吩咐的任务。后来，在狂怒中，他杀死了自己的幼子及妻子墨伽拉。最后，穿上妻子得伊阿尼拉送来的衣服，因溃疡而濒死，无法忍受痛苦，在埃塔山上为自己堆起柴堆，自焚而死。因此，尽管因他的卓越，他可能曾被相信是神，但鉴于这些事，他被相信是人了。\par

% structure: p0024 source=h2 target=section reason=relative-heading-rank; base=h2

\section{第八章——论埃斯库拉庇乌斯、阿波罗、玛尔斯、卡斯托尔和波鲁克斯，以及墨丘利和巴克斯。}

塔奎提乌斯（Tarquitius）记载，埃斯库拉庇乌斯（Æsculapius）出身可疑，因此被遗弃；后被猎人收养，由母犬之乳哺育，交给基戎（Chiron）教导。他住在埃皮达鲁斯（Epidaurus），葬于库诺苏莱（Cynosuræ），如西塞罗（Cicero）所说，被雷电击死。但他的父亲阿波罗（Apollo）并未不屑于照管他人的羊群以换取妻子；当他无意中杀死自己所爱的少年时，便将自己的哀悼题写在一朵花上。马尔斯（Mars）是最勇敢的人，却未能免于通奸的指控，因为他与奸妇一同被锁链捆绑，成为笑柄。\par

卡斯托尔（Castor）和波鲁克斯（Pollux）抢夺他人的新娘，但并非不受惩罚；荷马（Homer）以朴素真实而非诗意的信仰见证了他们的死亡与埋葬。墨丘利乌斯（Mercurius）与维纳斯（Venus）私通而生下安德洛吉努斯（Androgynus），他因发明了竖琴和角力场而应成为神祇。父神巴克斯（Bacchus）征服印度后，偶然来到克里特（Crete），看见忒修斯（Theseus）强迫并遗弃的阿里阿德涅（Ariadne）在岸边。于是他因爱而燃起激情，娶她为妻，并将她的冠冕，如诗人们所说，高悬于群星之间。众神之母本人，在丈夫被放逐并死后，居住在弗里吉亚（Phrygia）时，虽为寡妇且年迈，却爱慕一名美少年；因他不忠，她便将他阉割，使他变得柔弱；因此直到如今，她仍以加里（Galli）为司祭而喜悦。\par

% structure: p0027 source=h2 target=section reason=relative-heading-rank; base=h2

\section{第九章——论诸神可耻的行为。}

刻瑞斯（Ceres）除了因放荡行为，从哪里生出普洛塞庇娜（Proserpine）？拉托娜（Latona）除了因犯罪，从哪里生出她的双胞胎？维纳斯（Venus）曾屈从于神与人的情欲，当她在塞浦路斯（Cyprus）为王时，发明了娼妓行径，并命令妇女卖淫，以免唯独自己蒙羞。难道贞女密涅瓦（Minerva）和狄安娜（Diana）本身是贞洁的吗？那么埃里克托尼乌斯（Erichthonius）从哪里来？难道伏尔甘（Vulcan）将种子洒在地上，人便像菌类那样从那里生出？或者，为什么狄安娜将希波吕托斯（Hippolytus）放逐到偏僻之处，或将他交给一个女人，使他在陌生的树林中孤独度过一生，如今改名后被称为维尔比乌斯（Virbius）？这些事除了不洁之外，还意味着什么？而诗人们不敢承认。\par

% structure: p0029 source=h2 target=section reason=relative-heading-rank; base=h2

\section{第十章——论朱庇特及其放荡生活。}

但关于所有这些神的王和父——朱庇特（Jupiter），他们相信他掌有天上至高权能——他有何权能？他放逐了他的父亲萨图尔努斯（Saturnus）离开王国，并在其逃亡时持武器追赶。他有何节制？他纵情于各种情欲。因为他使伟人的妻子阿尔克墨涅（Alcmena）和勒达（Leda）因他的通奸而蒙羞；他也为一名少年的美貌所迷惑，在他狩猎并思考男子气概之事时，强行将他掳走，以便像对待女子那样对待他。我为何要提他玷污贞女的事？其数目之多，从他众子的数量可见一斑。唯独对特提斯（Thetis）他更为克制。因为曾有预言说，她所生的儿子将比其父更强大。因此他克制自己的爱欲，以免生出一个比他更伟大的人。可见他知道自己并非拥有完美的美德、伟大与权能，因为他害怕他自己对他父亲所做的事。那么，他为何被称为至善至大者？既然他以恶行玷污自己（这是不义与恶者所为），又惧怕比自己更强者（这是软弱与卑下者所为）。\par

% structure: p0031 source=h2 target=section reason=relative-heading-rank; base=h2

\section{第十一章——诗人们用以掩盖朱庇特丑行的各种象征。}

但有人会说，这些事是诗人们虚构的。诗人的习惯并非如此虚构，以至于你捏造全部，而是以形象、如同多彩的面纱掩盖事实本身。诗意的自由有此界限：不是可以凭空编造一切（那是虚假与愚昧者所为），而是可以合理改变某些事物。他们说朱庇特化为金雨以欺骗达娜厄（Danae）。金雨是什么？显然是金币；通过大量献上这些并倾入她的怀中，他以此贿赂腐蚀了她贞洁灵魂的软弱。同样，当他们想表示大量投枪时，也说到\Quote{铁雨}。他乘鹰劫走他的男宠。鹰是什么？实为军团，因为这种动物的形象是军团的徽帜。他骑牛渡海带走欧罗巴（Europa）。牛是什么？显然是一艘船，其守护神像被塑成公牛形状。因此，伊那科斯（Inachus）的女儿确实没有变成母牛，也没有以母牛之姿游水，而是乘着一艘有母牛形状的船逃脱了朱诺（Juno）的愤怒。最后，她被带到埃及后成为伊西斯（Isis），其航行在固定的一日被纪念，以追忆她的逃亡。\par

% structure: p0033 source=h2 target=section reason=relative-heading-rank; base=h2

\section{第十二章——诗人们并非虚构所有涉及诸神之事。}

你看，诗人们并非虚构一切，而是预表了某些事物，以便当他们说出真相时，能给他们称之为神者加上某种类似神性的东西；正如关于他们的王国也是如此。因为当他们说朱庇特通过抽签获得天空的王国时，他们要么意指奥林波斯山（Mount Olympus）——古代故事说萨图尔努斯（Saturnus）及后来的朱庇特曾居住于此——要么意指东方的一部分，那里似乎更高，因为光从那里升起；而西方区域较低，因此他们说普鲁托（Pluto）获得了下界；海洋给了尼普顿（Neptune），因为他拥有海岸与所有岛屿。许多事就这样被诗人们涂上了色彩；不知此理的人责备他们虚假，但只在言语上；实际上他们相信这些事，因为他们如此塑造神像，以至于当他们将神祇定为男女，并承认有些已婚、有些为父母、有些为子女时，他们显然同意诗人们；因为这些关系若无交合与生育便不能存在。\par

% structure: p0035 source=h2 target=section reason=relative-heading-rank; base=h2

\section{第十三章——历史学家欧赫墨罗斯所记载的朱庇特的行迹}

但让我们离开诗人；让我们转向历史，它既得到事实可靠性的支持，也得到时代古老性的支持。欧赫墨罗斯是麦西那人，一位非常古老的作者，他依据古老神庙的圣刻文字，记述了朱庇特的起源、他的功绩以及他所有的后裔；他还追溯了其他神灵的父母、他们的国度、行为、命令和死亡，甚至他们的坟墓。恩尼乌斯将这段历史译成了拉丁文，其原文如下：\par

\Quote{正如这些事所记载的，朱庇特及其兄弟的起源和亲族也是如此；圣书就是以这种方式传于我们的。}\par

因此，同一位欧赫墨罗斯记载说，朱庇特五次周游世界，将统治权分给朋友和亲属，为人类制定法律，并施行了许多其他恩惠之后，带着不朽的荣耀和永久的纪念，在克里特结束了生命，归于众神；他的坟墓在克里特的格诺索斯镇，上面用古希腊文刻着\OriginalTerm{匝恩克罗努}{\GreekText{Ζαν} \GreekText{Κρόνου}}，意思是\Quote{克罗诺斯之子朱庇特}。因此，从我所述之事可以清楚看出，他是一个人，曾在世上为王。\par

% structure: p0039 source=h2 target=section reason=relative-heading-rank; base=h2

\section{第十四章——历史学家所记载的萨图尔努斯和乌拉诺斯的行迹}

让我们转向更早的事，以便发现整个谬误的起源。据说萨图尔努斯是天神和地神所生。这显然难以置信；但如此记载自有其理由，不知此理者便会视之为神话。赫尔墨斯肯定乌拉诺斯是萨图尔努斯的父亲，圣史也如此教导。特里斯墨吉斯托斯曾说，完美学识之人极少，他在其中列举了他的亲属乌拉诺斯、萨图尔努斯和墨丘利。欧赫墨罗斯记载说，同一位乌拉诺斯是第一个在地上为王的人，用了这些话语：\Quote{起初，天神首先在地上拥有最高权力；他为自己与兄弟们建立并预备了那个王国。}\par

% structure: p0041 source=h2 target=section reason=relative-heading-rank; base=h2

\section{第二十章——论罗马特有的诸神}

我已经谈过各民族共有的宗教仪式。现在我要谈罗马人特有的诸神。谁不知道福斯图卢斯的妻子、罗慕路斯和瑞摩斯的乳母，为纪念她而设立了拉伦提纳利亚节，竟是一个娼妓？因此她被称作\OriginalTerm{卢帕}{Lupa}，并以野兽的形象呈现。同样，福拉和弗洛拉也是娼妓，前者是赫拉克勒斯的情妇（据维里乌斯记载）；后者以卖淫积攒了大量财富，立人民为继承人，为此举行弗洛拉利亚节以纪念她。\par

塔提乌斯将在主要下水道中发现的一尊妇女雕像奉为神圣，称之为克洛阿基娜女神。罗马人受高卢人围攻时，用妇女的头发制造投射武器；为此他们建造了祭台和神庙献给秃头维纳斯；又献给面包师朱庇特，因为他曾在梦中指点他们将所有谷物做成面包，投向敌人；此事之后，高卢人认为无法以饥饿降服罗马人，便弃围而去。图路斯·霍斯提略封恐惧和苍白为神。心智也受崇拜；但若他们有心智，我相信他们绝不会认为它应当受崇拜。玛尔凯路斯创立了荣誉和美德之神。\par

% structure: p0044 source=h2 target=section reason=relative-heading-rank; base=h2

\section{第二十一章——论罗马诸神的宗教仪式}

但元老院还设立了其他此类虚假的神祇——希望、信仰、和谐、和平、贞洁、孝爱；这一切本当真实地存在于人的心中，他们却错误地将它们放置在墙壁之内。尽管这些神祇在人之外没有实体存在，但我宁愿它们受崇拜，胜于枯萎或热病（这些不应当被奉为神圣，而应受咒诅）；胜于福尔纳克斯及其神圣的烤炉；胜于斯特克图斯（他最先教人用粪肥肥沃土地）；胜于穆塔女神（她生下了拉瑞斯）；胜于库米娜（她主管婴儿的摇篮）；胜于卡卡（她向赫拉克勒斯告发其牛群被盗，使他得以杀死自己的兄弟）。有多少其他怪诞可笑的虚构，谈及它们令人痛心！然而，我不愿忽略对泰尔米努斯的提及，因为据说他即使对朱庇特也不让步，尽管他只是一块未经雕刻的石头。他们以为他掌管边界，向他献上公开的祈祷，求他使卡皮托利山上的石头不动摇，并保存和扩展罗马帝国的疆界。\par

% structure: p0046 source=h2 target=section reason=relative-heading-rank; base=h2

\section{第二十二章——论法乌努斯和努马引入的宗教仪式}

法乌努斯是在拉丁姆首先引入这些愚行的人；他既向他的祖父萨图尔努斯献上血腥的祭品，又希望他的父亲皮库斯被当作神来崇拜，并将他的妻子兼姐妹法图阿·法乌娜置于众神之中，称她为善神。然后，在罗马，努马以新的迷信加重了那些野蛮和乡民的负担，设立了司祭职，将众神分配到家族和族群，以便使民众的凶猛精神从战争的追求中转移开。因此，卢基利乌斯嘲笑那些受虚妄迷信奴役者的愚行，引用了这些诗句：\par

\Quote{那些吓唬孩童的拉弥亚，即法乌努斯和努马·蓬皮利乌斯等人所设立的，对此他战栗；他将一切寄托于此。正如幼童相信每一尊青铜像都是活人，这些人也以为一切虚构之物都是真的：他们相信青铜像中有心。那是一座画廊；无物真实，一切皆是虚构。}\par

图利乌斯在论诸神之本性时也抱怨说，引入了虚假和虚构的神祇，由此产生了错误的意见、混乱的谬误和近乎老妪般的迷信；这一意见与其他意见相比应被视为更有分量，因为这些话出自一位兼为哲学家和司铎之口。\par

% structure: p0050 source=h2 target=section reason=relative-heading-rank; base=h2

\section{第二十三章——论蛮族的诸神和宗教仪式}

我们已经论及诸神；现在将论及他们神圣体制的礼仪与惯例。曾有人祭献给塞浦路斯的朱庇特，如特乌克尔所规定。同样，陶里人也将外乡人祭献给狄安娜；拉丁姆的朱庇特也以人血求得和解。在萨图尔努斯之前，依照阿波罗的神谕，六十岁的男人被从桥上扔进台伯河。而迦太基人不仅将婴儿祭献给同一萨图尔努斯，而且被西西里人击败后，为赎罪，他们祭献了二百名贵族之子。那些至今仍献给大神母和贝罗娜的祭献，也并不比这些更温和：在这些祭献中，司祭们不是用他人的血，而是用自己的血献祭；当他们自残身体时，他们不再是男人，却也未成为女人；或者切割肩膀，用鲜血洒满可憎的祭台。但这些是残忍的。\par

让我们来看那些温和的。伊西斯的圣礼仅仅显示她如何失去又找到她的小儿子，即奥西里斯。起初，她的司祭和侍从剃光全身毛发，捶胸哭嚎、寻找，模仿他母亲当时的情绪；随后男孩被狗头神找到。于是悲哀的礼仪以欢喜结束。刻瑞斯的奥义也类似：点燃火炬，整夜寻找普洛塞庇娜；找到后，整个仪式以祝贺和投掷火炬结束。兰普萨库斯人向普里阿普斯献上一头驴作为适当的祭品。林都斯是罗得岛的一个城镇，那里举行赫拉克勒斯的圣礼，其间伴有辱骂。因为赫拉克勒斯从农夫那里夺走并杀死了牛，农夫以辱骂报复；后来赫拉克勒斯自己被任命为司祭，规定他本人和后来的司祭都要以同样的辱骂举行祭献。但克里特朱庇特的奥义显示他是如何被藏匿或养育的。山羊在他身旁，阿玛尔特亚以此喂哺男孩。众神之母的圣礼也显示同样的事。因为当时科里班忒斯以头盔的叮当声和盾牌的敲击声掩盖男孩的哭声，如今在圣礼中重现这一情景；但以钹代替头盔，以鼓代替盾牌，使萨图尔努斯听不见男孩的哭声。\par

% structure: p0053 source=h2 target=section reason=relative-heading-rank; base=h2

\section{第24章 论圣礼与迷信的起源}

这些是诸神的奥义。现在让我们也探询迷信的起源，以便查明它们由谁、在何时建立。狄迪莫斯在题名为\WorkTitle{品达注释}的书中说，梅利塞乌斯是克里特人的王，他的女儿是阿玛尔特亚和梅利莎，她们用羊奶和蜂蜜喂养朱庇特；他引入了新的礼仪和仪式，并且是第一个向诸神献祭的人，即向维斯塔（又名特勒斯）献祭——为此诗人写道：\par

\Quote{诸神中首先的，\newline{}特勒斯，}\par

随后向众神之母献祭。但欧赫墨鲁斯在其\WorkTitle{圣史}中说，朱庇特本人取得统治权后，在许多地方为自己建造了庙宇。因为当他巡游世界时，每到一处，他便与当地首领结为朋友，享有宾客之权；为保存这一记忆，他下令为他建造庙宇，并由与他结盟的人每年举行节庆。如此，他使对自己的崇拜遍及所有土地。他们生活的时代容易推断。因为塔鲁斯在其史书中写道，亚述王贝鲁斯（巴比伦人所崇拜者），是萨图尔努斯的同代人和朋友，比特洛伊战争早三百二十二年；而特洛伊陷落至今已一千四百七十年。由此显然，从人类因建立新的崇拜形式而堕入错误，至今不过一千八百年。\par

% structure: p0057 source=h2 target=section reason=relative-heading-rank; base=h2

\section{第25章 论黄金时代、神像及普罗米修斯（他首先塑造了人）}

因此，诗人们很有理由地说，萨图尔努斯统治时的黄金时代被改变了。因为那时无人崇拜诸神，他们只认识独一天主。后来，他们屈从于脆弱和尘世之物，崇拜木、铜、石的偶像，于是从黄金时代转变为铁时代。因为他们失去对天主的认识，并断绝人类社会的那唯一纽带，开始互相侵扰、掠夺和征服。但如果他们举目仰望天主——祂将他们提升到天上，并使他们仰望祂自己及祂本身——他们就不会俯伏跪拜尘世之物了。卢克莱修严厉斥责这种愚昧，说：\par

\Quote{他们因对诸神的惧怕而卑抑灵魂，使之重压在地。}\par

因此他们战栗，却不明白畏惧你所制造的东西是何等愚蠢，或指望从那些无言无觉、既看不见也听不见祈求者的东西那里获得任何保护是何等愚昧。那么，它们能有什么威严或神性呢？它们曾处于人的权下，能否不被制造，或被制成别的东西，甚至现在仍是如此？因为它们易受损伤，可能被偷走，若非受到法律和人的守护。那么，向这样的神祇献上最精选的祭品、供奉礼物、献上昂贵的衣服，仿佛那些不动的东西能用它们似的，这样的人看来神志清醒吗？所以，西西里的暴君狄奥尼修斯在征服希腊后劫掠并嘲弄希腊的神祇，是有理由的；在他犯下这些亵渎行为之后，他顺风顺水返回西西里，统治王国直到老年；而受损的神祇竟无法惩罚他。\par

轻视虚妄、归向天主、持守你从天主所领受的处境、持守你的名字，岂不更好！因为人之所以被称为\OriginalTerm{人}{\GreekText{ἄνθρωπος}}，正是因为他向上仰望。然而向上仰望的人，是仰望那居于高天的真实永生天主；他不仅以知觉和理智，也以高举的目光和面容寻求其灵魂的造主与父。但若有人使自己服从于尘世卑微之物，显然是把低于自己的东西置于自身之上。因为人本身是天主的化工，而偶像是人的化工，人的化工不能优于神的化工；正如天主是人的父，人也是雕像的父。因此，崇拜自己所制造之物的人是愚蠢无知的。这种可憎而愚蠢的手艺的创始者是普罗米修斯，他出生于朱庇特叔父伊阿珀托斯。因为当初朱庇特获得至高的统治权后，想把自己立为神并建立庙宇，正寻找能模仿人形的人，那时普罗米修斯活着，他用稠泥塑造了人像，其逼真程度使艺术的创新和精巧令人惊叹。最后，他同时代的人以及后来的诗人，都将他传颂为真正活人的制造者；而我们每当赞美雕像时，就说它们活着、呼吸。他确实是陶像的发明者。但后辈追随他，既用大理石雕刻，也用青铜铸造；随后随着时间的推移，又加上了金和象牙的装饰，以致不仅形象，连光辉本身也能炫人眼目。就这样，被美丽所迷惑，忘记了真正的威严，有感觉的生物认为无感觉之物、有理性的生物认为无理性之物、有生命的生物认为无生命之物应当受他们崇拜和尊敬。\par

% structure: p0062 source=h2 target=section reason=relative-heading-rank; base=h2

\section{第二十六章——论崇拜元素与星辰。}

现在让我们也驳斥那些视世界元素为神的人，即视天、太阳、月亮为神的人；因为他们不认识这些的造主，反而赞美并崇拜这些受造物本身。这种错误不仅属于无知之人，也属于哲学家；因为斯多亚派认为所有天体都应算在众神之列，因为它们都有固定而有规律的运行，从而最恒定地保持着相继季节的变化。那么，它们并不具有自主的运动，因为它们遵从规定好的法则，且显然不是出于自身的知觉，而是出于至高造主的化工，祂如此安排它们，使它们完成无误的轨道和固定的循环，藉以交替昼夜、冬夏。但是，若人们赞美这些结果，赞美它们的运行、光辉、规律、美丽，他们就应当明白，比这些更美丽、更辉煌、更强大的是制造者和设计者本身，即天主。然而，他们以人眼所见之物来估量神性，不知道可见之物不能永恒，而永恒之物不能为凡眼所辨识。\par

% structure: p0064 source=h2 target=section reason=relative-heading-rank; base=h2

\section{第二十七章——论人的受造、罪恶与惩罚；以及善恶天使。}

还剩最后一个主题：既然据历史记载，众神常通过预兆、梦境、神谕以及惩罚亵渎者来显示其威严，那么我要阐明是什么原因产生了这种效果，以使现在也没有人陷入古人曾陷入的同样陷阱。当天主按照祂卓越的威严从无中创造了世界，用光体装饰了天空，用生物充满了地和海之后，祂用泥土造了人，按照自己的肖像塑造了他，向他吹气使他成活，把他安置在祂所种植的果园中，园中有各种果树；祂命他不得吃一棵树上的果子，这棵树有分辨善恶的知识，并警告他，若如此行，他将丧失生命，但若遵守天主的诫命，他将保持不死。然后那蛇，原是事奉天主者之一，嫉妒人因被造而能不朽，就用诡计引诱人违犯天主的诫命和法律。这样，人确实得到了善恶的知识，却丧失了天主所赐予他的永生。\par

因此，祂将罪人逐出圣地，流放到这个世界，使他必须劳苦才能谋生，按其所应得遭受艰难与困苦；祂又用火墙环绕那园子，使得在审判之日以前，无人能秘密进入那永福之地。随后，按照天主的判决，死亡临到人身上；然而，人的生命虽已变为暂时，其界限仍为一千年，直到洪水时代，人类寿命的长度仍是如此。因为洪水之后，人的寿命逐渐缩短，减少到一百二十年。但那蛇——因其行为而得名为魔鬼，即控告者或告密者——并未停止迫害他从起初所欺骗的人的后裔。最终，在嫉妒的冲动下，他怂恿世上首先出生的人杀害自己的兄弟，为要消灭头生的两人中的一人，并使另一人成为弑亲者。此后，他也不停止在每一代人中向人的胸怀注入恶意的毒液，败坏和堕落他们；总之，使他们陷入如此深重的罪恶，以致正义的事例如今罕见，人却如同野兽般生活。\par

但天主看见这事，便派遣祂的天使去教导人类，保护他们远离一切邪恶。祂命令他们远离属世的事物，免得被任何玷污所污染，失去天使的尊荣。然而，那狡猾的控告者，当他们在人间逗留时，也引诱他们陷入享乐，使他们与女子玷污自己。于是，他们被天主的判决定罪，因自己的罪恶被逐出，丧失了天使的名称和实质。这样，他们成了魔鬼的仆役，为了给自己的堕落寻求安慰，便致力于毁灭那些他们本为保护而来的人。\par

% structure: p0068 source=h2 target=section reason=relative-heading-rank; base=h2

\section{第二十八章——论魔鬼及其邪恶作为。}

这些就是魔鬼，诗人们常在诗篇中提及，\Person{赫西俄德}称他们为人的监护者。因他们用诱惑和诡计如此说服人，以致人相信他们就是神。最后，\Person{苏格拉底}宣称，从幼年起就有一个魔鬼作为他生命的监护者和指引者，没有它的同意和命令，他什么也不能做。因此，他们依附于个人，以\OriginalTerm{家神}{Genii}或\OriginalTerm{家神}{Penates}的名义占据房屋。为他们建造庙宇，每日向他们奠酒如同向\OriginalTerm{家神}{Lares}致敬，尊他们为避祸者。从起初，为使人远离对真天主的认识，他们引入了新的迷信和众神的崇拜。他们教导说，应将已故君王的记忆圣化，为他们建造庙宇，雕刻偶像——并非为减少天主的荣耀或增加他们自己的荣耀（他们因犯罪已失去荣耀），而是为从人那里夺走生命，剥夺他们对真光的希望，免得人达到他们自己所失落的不朽天赏。他们又揭示了占星术、占卜和预言术；尽管这些事情本身是虚假的，但作为邪恶的创造者，他们亲自掌管和调整这些事，使人相信它们是真实的。他们还发明了巫术的把戏，欺骗人的眼睛。借着他们的帮助，使有者看似无，无者看似有。他们自己发明了召魂术、神谕和预言，用模棱两可的结果的虚假占卜来迷惑人的心智。他们出现在庙宇和一切祭献中；通过展示一些欺骗性的奇事，使在场的人惊讶，从而欺骗人，使人相信在偶像和雕像中有神力存在。他们甚至作为轻盈的灵体秘密进入身体，在受损的肢体中引发疾病，当用祭献和誓愿平息他们时，他们又能移除。他们送出的梦充满恐惧，好使人呼求他们，或者梦的应验与事实相符，以增加人对他们的敬拜。有时他们也针对亵圣者施行某种报应，使看见的人更加胆怯和迷信。就这样，他们用欺诈给人类蒙上黑暗，使真理受压制，使至高无上、无可比拟的天主之名被遗忘。\par

% structure: p0070 source=h2 target=section reason=relative-heading-rank; base=h2

\section{第二十九章——论天主的忍耐与眷顾。}

但有人会说：那么，真天主为什么允许这些事发生？祂为什么不除掉或消灭恶人？实际上，祂为什么从起初就赋予魔鬼能力，使存在一个可以败坏并毁灭一切者？我将简要说明祂为何愿意如此。我问：美德是善还是恶？不能否认它是善。如果美德是善，那么相反地，恶习就是恶。如果恶习之所以为恶，是因为它反对美德，而美德之所以为善，是因为它战胜恶习，那么，没有恶习便不可能有美德；如果你除去恶习，美德的功绩也会被除去。因为没有敌人，就不能有胜利。因此，没有恶，善便不能存在。\par

\Person{Chrysippus}，一位心智活跃的人，在讨论天主眷顾时看到了这一点，并指责那些认为善由天主造成，却说恶并非如此造成的人是愚昧的。\Person{Aulus Gellius}在他的《\WorkTitle{阿提卡之夜}》一书中解释了他的观点，这样说道：\Quote{那些不认为世界是为天主和人的缘故而造、且人间事务由天主眷顾管理的人，以为他们使用了一个有力的论据，当他们这样说：如果真有天主眷顾，就不会有邪恶。因为他们说，没有什么比这更不符合天主眷顾：在这个据说天主为人而造的世界上，烦恼和邪恶的力量竟如此巨大。针对这些，\Person{Chrysippus}在论天主眷顾的第四卷书中辩论时说：再没有比那些认为好事能够存在、假如同地没有邪恶的人更愚蠢的了。因为既然善与恶相对，它们必然相互对立，并仿佛建立在相互反对的支撑上。因此，没有一个对立面缺少另一个对立面。因为如果没有不义，怎能感知正义？除了消除不义，正义又是什么？同样，刚毅的性质，除非将怯懦置于其旁，否则无法理解；节制的性质，除非有放纵；明智的性质，除非有不明智。按照同样的原理，他说，为什么愚昧之人不也要求这一点：应当有真理而没有虚假呢？因为善与恶、顺境与逆境、快乐与痛苦是并存的。因为两者在相反的两极上彼此相连，正如\Person{Plato}所说，如果你拿走一个，你就拿走了另一个。}因此，你看到了我常说的：善与恶如此相连，以致一方不能没有另一方。所以天主以最大的远见，将德行的素材置于祂为此目的而造的恶之中，好为我们设立一场竞赛，在其中祂要赏赐胜利者永生的赏报。\par

% structure: p0073 source=h2 target=section reason=relative-heading-rank; base=h2

\section{第三十章　论假智慧}

我想我已经教导过，献给诸神的崇敬不仅是邪恶的，而且也是徒劳的，无论是由于他们是死后被纪念的人，还是由于像本身没有感觉和听闻，因为它们由泥土构成，而人本应当仰望天上之事，却不应屈服于地上之物；或由于要求这些宗教崇拜的灵是不圣洁和不洁的，因此被天主的判决所定罪，堕落到地上，而如果你希望作真天主的追随者，就不应当服从那些你本就高过其权能者的权力。既然我们已经谈论过假宗教，接下来也应当讨论假智慧的课题，就是哲学家所宣称的——这些人富有极大的学识和口才，却远离真理，因为他们既不认识天主，也不认识天主的智慧。虽然他们聪明而有学问，但因为他们的智慧是属人的，我将不惧怕与他们争辩，好使虚假容易为真理所胜、属地之事容易为属天之事所胜得以显明。\par

他们这样定义哲学的性质：哲学是对智慧的爱或追求。因此它不是智慧本身；因为爱者必须与被爱者不同。如果是追求智慧，即使如此哲学也不等于智慧。因为智慧是所寻求的对象本身，而追求是寻求它的行为。因此这个词的定义或含义清楚地表明哲学不是智慧本身。我要说它甚至也不是追求智慧，因为智慧并不包含在这追求之中。因为谁能说他专心追求某种他根本无法达到的东西呢？专心于医学、文法或修辞的人，可以说他勤于学习那门技艺；但当他学会后，他就被称为医生、文法家或演说家。同样，那些勤于追求智慧的人，在学会之后，本应被称为智者。但由于他们一生被称为智慧的追求者，很明显那并不是追求，因为不可能达到追求中所寻求的对象本身，除非那些追求智慧直到生命尽头的人，在另一个世界将成为智者。现在，每个追求都有其目的。因此，没有目的的追求就不是正确的追求。\par

% structure: p0076 source=h2 target=section reason=relative-heading-rank; base=h2

\section{第三十一章　论知识与推测}

此外，有两个东西似乎属于哲学的主题——知识和假设；如果它们被取消，哲学就会完全崩溃。但哲学家中的首要人物已经把两者都从哲学中取消了。\Person{苏格拉底}取消了知识，\Person{芝诺}取消了假设。我们来看看他们这样做是否正确。智慧，按照\Person{西塞罗}的定义，是关于神圣事物和人类事物的知识。如果这个定义是正确的，那么智慧就不在人的能力范围内。因为哪个凡人能自认为知道神圣和人类的事物呢？我对人类事物不置一词；尽管它们与神圣事物相关，但由于它们属于人，我们可以承认人可能知道它们。但人靠自己不能知道神圣事物，因为他是人；而知道神圣事物的人必须是神圣的，因此就是天主。但人既不是神圣的，也不是天主。因此，人不能靠自己完全知道神圣事物。所以，除了天主，或者那个被天主教导的人，没有人是智慧的。但他们，因为他们既不是神，也没有被天主教导，就不能是智慧的，即通晓神圣和人类事物。因此，\Person{苏格拉底}和学园派取消知识是正确的。假设也不适合智慧之人。因为每个人都假设他所不知道的东西。而假设你知道你不知道的东西，是鲁莽和愚蠢。因此，\Person{芝诺}取消假设是正确的。所以，如果人里面没有知识，也不应该有假设，那么哲学就被连根拔起了。\par

% structure: p0078 source=h2 target=section reason=relative-heading-rank; base=h2

\section{第32章——论哲学家的学派及其分歧。}

此外，哲学不是统一的；它被分为许多学派，散播到众多且不一致的意见中，没有固定的状态。因为他们都各自攻击和骚扰彼此，没有哪一个不在其他派的判断中被谴责为愚蠢，既然成员之间明显彼此不和，整个哲学体系就走向毁灭。因此，后来产生了学园派。当那个学派的领袖们看到哲学被相互反对的哲学家们彻底推翻时，他们向所有人开战，以摧毁所有哲学家的所有论证；而他们自己除了断言一件事——什么都不能知道——之外，什么也不断言。这样，通过取消知识，他们推翻了古代哲学。但他们自己甚至没有保留哲学家的称号，因为他们承认自己的无知，因为对所有事情无知不仅不是哲学家应有的，甚至不是人应有的。因此，哲学家们因为没有辩护，必须用彼此造成的伤口互相毁灭，哲学本身也必须用其自己的武器完全消耗并终结自己。但有人说，只有自然哲学这样让位。道德哲学呢？它建立在任何牢固的基础上吗？让我们看看哲学家们是否在这个至少涉及生活状况的领域里一致。\par

% structure: p0080 source=h2 target=section reason=relative-heading-rank; base=h2

\section{第33章——生命中应追求的至善是什么。}

什么是至善是必须探究的对象，以便我们整个生活和行动可以指向它。当探究人的至善时，它应当被确定为这样的：首先，它只与人有关；其次，它特别属于心灵；最后，它要通过德性追求。因此，我们来看看哲学家所指出的至善是否是这样的：它既不涉及不会说话的动物，也不涉及身体，并且没有德性就无法获得。\par

昔兰尼学派的创始人\Person{亚里斯提卜}，认为身体的快乐是至善，他应当被从哲学家的行列和人类社会中移除，因为他把自己比作野兽。\Person{希罗尼穆斯}的至善是无痛苦，\Person{狄奥多鲁斯}的至善是停止痛苦。但其他动物避免痛苦；当它们没有痛苦或停止痛苦时，它们就高兴。那么，如果人的至善被判断为与兽类相同，人还有什么区别？\Person{芝诺}认为至善是合乎自然地生活。但这个定义是普遍的。因为所有动物都合乎自然地生活，每种动物都有自己的本性。\par

\Person{伊壁鸠鲁}主张至善是灵魂的快乐。什么是灵魂的快乐，除了喜悦，灵魂在其中大多放纵自己，或转向嬉戏或大笑？但这种善甚至发生在不会说话的动物身上，当它们吃饱草料后，就放松到欢乐和放纵中。\Person{狄诺玛库斯}和\Person{卡利丰}赞成可敬的快乐；但他们要么和\Person{伊壁鸠鲁}说的一样，认为身体的快乐是可耻的；要么如果他们考虑身体的快乐部分是低贱的部分是可敬的，那么这就不是归于身体的至善。逍遥学派将至善归结为灵魂、身体和命运的善。灵魂的善可以赞同；但如果它们需要帮助才能完成幸福，那么它们显然软弱。但身体和命运的善不在人的能力范围内；而且现在至善不是被归于身体或置于我们之外的事物，因为这种双重善甚至延伸到牲畜，它们需要身体健康和适当的食物。斯多葛学派被认为持有更好的观点，他们说德性就是至善。但德性不能是至善，因为如果它是忍受邪恶和劳苦，它本身并不幸福；它应该产生和造成至善，因为没有最大的困难和劳动就无法获得至善。但事实上，\Person{亚里士多德}远远偏离了理性，他将荣誉与德性联系起来，好像德性任何时候都能与荣誉分离，或与卑贱结合一样。\par

皮浪主义者\Person{赫里鲁斯}以知识为至善。这确实属于人，且只属于灵魂，但它可以在没有德行的情况下发生。因为一个人若只是听闻某事，或略加阅读便获得知识，并不能被视为幸福；这也不是至善的定义，因为知识可能关乎恶事，或至少是无用之事。即使你以辛劳获得了善恶之事的正确知识，它也不是至善，因为知识并非为自身而被寻求，而是为了别的缘故。因为学习技艺，是为了它们能成为我们谋生、荣誉甚至享乐的手段；显然，这些不可能是至善。因此，哲学家们即使在道德哲学中也不遵循规则，因为他们就在那塑造生命的关键讨论上彼此相异。因为当一些人引导人走向享乐，另一些人走向荣誉，还有一些人走向本性，另有一些人走向知识；一些人追求财富，另一些人回避财富；一些人完全无感于痛苦，另一些人忍受苦难时，这些诫命不可能相等或相似。在所有这一切中，正如我之前所表明的，他们背离了理性，因为他们不认识天主。\par

% structure: p0085 source=h2 target=section reason=relative-heading-rank; base=h2

\section{第三十四章——论人生来为正义}

现在让我们看看为智者所提议的至善是什么。人生来为正义，不仅圣经教导了这一点，有时连这些哲学家自己也承认。因此，\Person{西塞罗}说：\Quote{但在学者们讨论的所有事情中，最卓越的无疑是清楚理解我们生来为正义。}这极为真实。因为我们不是生来为邪恶的，我们是社会性和群居的动物。野兽被造是为了施展其凶猛；因为除了靠掠夺和流血之外，它们无法生存。然而，即使它们被极度饥饿所迫，仍会克制自己，不伤害同类。鸟类也是如此，它们必须靠其他动物的尸体为食。那么，人与人通过语言交流和情感相通而结合，人应当同情并爱另一个人，这岂不更合宜？因为这就是正义。\par

但由于智慧唯独赐予人，使他能认识天主，而这一点就构成了人与哑巴动物之间的区别，正义本身便包含两种义务。一种是对天主如同对父，另一种是对人如同对兄弟；因为我们都由同一天主所造。因此，有一个公正而恰当的说法：智慧是关于神圣与人类事务的知识。因为我们应该知道我们欠天主什么，欠人什么；即对天主有宗教，对人行仁爱。前者属于智慧，后者属于德行；而正义包含两者。因此，如果显然人生来为正义，那么义人就必须受苦受难，以便他操练所赋予的德行。因为德行就是忍受苦难。他会视享乐为恶而回避；他会轻视财富，因为财富是脆弱的；如果他拥有财富，他会慷慨施舍以救助困苦者；他不会渴求荣誉，因为荣誉短暂而易逝；他不伤害任何人；若受伤害，他不报复；也不向掠夺他财产的人复仇。因为他认为伤害他人是不合法的；若有人强迫他背离天主，他绝不拒绝酷刑或死亡。如此，他必然生活在贫穷、卑微、侮辱甚至折磨之中。\par

% structure: p0088 source=h2 target=section reason=relative-heading-rank; base=h2

\section{第三十五章——论不朽为至善}

那么，如果正义与德行在生活中只有苦难，它们的好处何在呢？但如果德行轻视一切地上的事物，最智慧地忍受一切苦难，甚至在尽本分时承受死亡本身，却不可能没有赏报，那么剩下来的唯有不朽是它的赏报。因为如果哲学家们愿意，幸福生活是人的分，而在这一点上他们并无分歧，那么不朽也属于他。因为唯有不腐朽的才是幸福的；唯有永恒的才是不腐朽的。因此不朽是至善，因为它既属于人，也属于灵魂，也属于德行。我们被引向这个目标；我们生来就是为了达到它。因此天主向我们提出德行和正义，为使我们因劳苦而获得那永恒的赏报。关于不朽本身，我们将在适当的地方谈论。剩下的逻辑哲学对幸福生活毫无贡献。因为智慧不在于言辞的安排，而在于内心与情感。但如果自然哲学是多余的，逻辑哲学也是多余的，而哲学家们在唯一必要的道德哲学中犯了错误，因为他们无论如何都找不到至善，那么全部哲学都是空洞无用的，因为它既不能理解人的本性，也不能履行其本分与职责。\par

% structure: p0090 source=h2 target=section reason=relative-heading-rank; base=h2

\section{第三十六章——论哲学家——即\Person{伊壁鸠鲁}和\Person{毕达哥拉斯}}

我已简略论述了哲学，现在也要略谈哲学家。这是 \Person{伊壁鸠鲁} 的特别学说：没有天主眷顾。同时他并不否认神的存在。在两方面他都违背理性。因为若有神存在，则必有天主眷顾。否则我们无法形成任何对天主的可理解概念，因为预见正是祂的特有领域。但 \Person{伊壁鸠鲁} 说祂毫不关心任何事物。因此祂不仅不顾人类事务，也不顾天上之事。那么，你如何或从何断言祂存在？因为当你否定了神圣眷顾与关怀，自然会得出你应完全否认天主存在的结论；而你如今在名义上保留了祂，实际上却摒弃了祂。那么，世界从何起源，若天主毫不关心任何事物？他说，有微小的原子，既不可见亦不可触，万物由这些原子的偶然碰撞产生，并不断生成。若它们无法被身体任何部分看见或感知，你何以知道它们存在？其次，若它们存在，它们以何种心智聚合以产生任何效果？若它们是光滑的，则无法结合；若它们有钩和角，则可分割；因为钩和角突出，可被切断。但这些都毫无意义且无益。为何我还要提及他甚至主张灵魂可朽？反驳他的不仅有所有哲学家和普遍信念，还有诗人的回答、西卜拉女祭司的预言，以及最后先知们自身的神圣声音；以致令人惊奇的是，唯独 \Person{伊壁鸠鲁} 竟将人的处境置于与牲畜和野兽同等的水平。\par

\Person{毕达哥拉斯} 又如何？他是第一位被称为哲学家的人，他判定灵魂诚然不朽，但会转入其他身体——或牛、或鸟、或野兽。与其这样被罚转入其他动物的身体，岂不与身体一同消亡更好？与其曾经具有人形而后活为猪或狗，岂不根本不存在更好？这位愚人，为使他的说法可信，自称在 \Place{特洛伊} 战争时曾是 \Person{欧福耳玻斯}，被杀后进入其他动物形体，最终成为 \Person{毕达哥拉斯}。幸福的人！唯独他被赐予如此巨大的记忆；或者不如说不幸，当他变成一只羊时，仍被允许不忘记自己是什么！但愿只有他一人如此愚昧！他还找到了一些相信他的人，甚至包括一些有学问的人，愚昧的遗产传给了他们。\par

% structure: p0093 source=h2 target=section reason=relative-heading-rank; base=h2

\section{第三十七章——论 \Person{苏格拉底} 及其矛盾}

在他之后，\Person{苏格拉底} 在哲学中居首位，甚至被神谕宣称为最有智慧，因为他承认自己只知道一件事——即 \Quote{他一无所知}。基于这神谕的权威，自然哲学家应当自制，以免探究他们无法知道的事，或以为知道他们并不知道的事。然而，让我们看看 \Person{苏格拉底} 是否如 \Person{皮提亚神} 所宣告的那样最有智慧。他常使用这句谚语：\Quote{我们之上的事与我们无关}。他现在已超越了他的见解界限。因为那个说他只知道一件事的人，找到了另一件事可谈，仿佛他知道它；但那只是徒劳。因为天主，显然在我们之上，应被寻求；宗教应被接受，它唯独将我们与禽兽区分开来，而 \Person{苏格拉底} 不仅拒绝，甚至嘲笑宗教，他指着鹅和狗发誓，仿佛他真不能指着 \Person{埃斯库拉庇俄斯} 发誓，他曾向祂许愿献一只公鸡。看这智慧人的祭献！因为他自己无法在临终时奉献此祭，他请求朋友在他死后代为还愿，以免他作为债务人被拘留在阴间。他确实既声明自己一无所知，也证实了此言。\par

% structure: p0095 source=h2 target=section reason=relative-heading-rank; base=h2

\section{第三十八章——论 \Person{柏拉图}，其学说更接近真理}

他的弟子 \Person{柏拉图}——\Person{图利乌斯} 称他为哲学家之神——独自在所有哲学研究中如此接近真理；然而，由于他不认识天主，他在许多事上如此失败，以致无人陷入更严重的错误，特别是因为他在 \WorkTitle{理想国} 一书中希望一切事物为所有人共有。这若涉及财产，虽不公正，尚可容忍。因为若有人因自身勤勉而拥有多于他人，不应损害任何人；若有人因自身过错而拥有较少，也不得受益。但如我所说，这在某种程度上尚可容忍。妻子和子女也要公有吗？血统没有分别、种族无从确定？没有家族、没有亲属、没有姻亲，一切都混乱无序，如同畜群？男人没有自制，女人没有贞洁？若双方之间没有确定或专一的爱，哪有什么夫妻之情？谁会对父亲恭敬，若不知自己从何而生？谁会爱儿子，若认为他不是自己亲生的？此外，他向妇女开放元老院，将战争、官职和指挥权委托给她们。但若妇女执行男人的职责，那城邦将遭受何等灾难！不过这一点以后有机会再详论。\par

斯多葛派的宗师\Person{芝诺}赞扬德行，却认为怜悯——一种极大的德行——应当被切除，仿佛它是一种心灵的疾病；然而怜悯同时为天主所珍视，也为人类所必需。因为，有谁在陷入任何灾祸时，不愿意被人怜悯，并渴望那些能援助他们的人伸出援手呢？而这种援手若不是出于怜悯之情，是不会被唤起以提供帮助的。尽管他把这称为人道与虔敬，但他并没有改变事情本身，只是改了名称。这种情感是唯独赐予人的，使我们能通过相互帮助来减轻自己的软弱；谁若除去这种情感，就是使我们沦入野兽的生活。因为他声称一切过错都是相等的，这种说法出自于那种不人道，他也以此攻击怜悯为一种疾病。因为一个不分过错轻重的人，要么认为轻罪应受重罚——这是残忍的法官所为；要么认为重罪应受轻罚——这是无价值的法官所为。无论哪种情况，都对国家有害。因为如果最大的罪行被轻轻放过，恶人的胆量就会增长并做出更大胆的事；而如果对轻罪施加过于严厉的惩罚，既然没有人能免于过错，许多本可通过改正而变得更好的公民就会陷入危险。\par

% structure: p0098 source=h2 target=section reason=relative-heading-rank; base=h2

\section{第三十九章——论各种哲学家与对跖人}

这些事情确实微不足道，但却出自同样的谬误。\Person{色诺芬尼}说月球球体是我们地球的十八倍大，并且在其范围之内包含着另一个地球，居住着各种人类和动物。关于对跖人，人们既不能听也不能说而不发笑。有人郑重其事地断言，我们应当相信存在着双脚与我们相对的人。\Person{阿那克萨戈拉}的胡言乱语更可容忍，他说雪是黑色的。不仅他们的言论，而且他们的一些行为也是可笑的。\Person{德谟克利特}荒废了他父亲留给他的土地，任其成为公共牧场。\Person{第欧根尼}和他的犬儒学派之徒，声称在轻视万物中拥有伟大而完美的德行，却宁愿乞讨为生，也不愿通过诚实的劳动或拥有任何财产来谋生。无疑，一位智者的生活理应成为他人生活的榜样。如果人人都效仿这些人的智慧，国家将如何存在？但也许这些犬儒学派的人能提供谦逊的榜样，他们与妻子在公共场合生活。我不知道他们如何能捍卫德行，既然他们已夺去了谦逊。\par

比他们好不了多少的是\Person{亚里斯提卜}，我相信，为了取悦他的情妇\Person{莱依斯}，他创立了昔兰尼学派，把至善归于身体快乐，这样他的过错就不缺乏权威，他的恶行就不缺乏学识。那些为了被称为藐视死亡而自杀的人，是否更应被视为更有毅力？\Person{芝诺}、\Person{恩培多克勒}、\Person{克吕西波}、\Person{克里安提斯}、\Person{德谟克利特}和效仿他们的\Person{加图}，不知道按照神权与神法，自杀的人犯了谋杀罪。因为正是天主将我们安置在这血肉的居所；正是祂赐给我们这身体的临时住所，要我们按祂的意愿居住多久就住多久。因此，未经天主命令而擅自离开，应被视为不虔敬。所以，不可对自然界施加暴力。祂知道如何摧毁自己的工程。若有人对那工程施以不虔敬之手，撕断神圣工艺的纽带，他就是企图逃离天主；无论生死，无人能逃脱祂的审判。因此，我上面提到的那些人是可诅咒的、不虔敬的，他们甚至教导了自愿死亡的理由，好像他们做自杀者还不够有罪，还要教唆别人也行此恶事。\par

% structure: p0101 source=h2 target=section reason=relative-heading-rank; base=h2

\section{第四十章——论哲学家的愚妄}

哲学家们有数不清的言行可以显明他们的愚妄。因此，既然我们无法一一列举，几例就够了。只要知道哲学家既不是正义的老师——因为他们对正义一无所知——也不是德行的老师——因为他们虚假地夸耀德行，这便足够了。因为他们能教导什么呢？他们常常承认自己的无知。我不提\Person{苏格拉底}，他的观点人所共知。\Person{阿那克萨戈拉}宣称万物都笼罩在黑暗之中。\Person{恩培多克勒}说感官寻求真理的道路是狭窄的。\Person{德谟克利特}断言真理深埋在一口深井中。因为他们哪里都找不到真理，所以他们断言至今尚无智者存在。既然人的智慧不存在（苏格拉底在\Person{柏拉图}的著作中说），我们就当追随那神圣的智慧，并感谢天主，因祂已将之启示并交付给我们。我们当为自己庆贺，因着天主的恩赐，我们拥有了真理和智慧，那是许多心灵经过如此多世纪所探求、而哲学却未能发现的。\par

% structure: p0103 source=h2 target=section reason=relative-heading-rank; base=h2

\section{第四十一章——论真宗教与智慧}

如今，我们既已驳斥了虚假的宗教（即关于敬拜诸神的宗教）与虚假的智慧（即关于哲学家的智慧），就让我们来谈真实的宗教与真实的智慧。而且，我们必须同时谈论它们两者，因为它们紧密相连。因为敬拜真天主，这正是智慧，而非其他。那位至高无上、万物的创造者，祂按自己的肖像造了人，为此唯独赐予人理性的恩赐，好使人能以孝爱之情报答祂——祂是人的圣父与主——并借由履行这种虔敬与服从，赢得不朽生命的赏报。这是一个真实而神圣的奥迹。但在那些虚假之中，却没有一致之处。哲学中并不举行神圣礼仪，神圣事物中也不探讨哲学；因此，他们的宗教是虚假的，因为它不具有智慧；他们的智慧也是虚假的，因为它不具有宗教。但两者相结合之处，必然有真理；如此，若问真理本身是什么，便可正当地称之为\Quote{智慧的宗教}或\Quote{宗教的智慧}。\par

% structure: p0105 source=h2 target=section reason=relative-heading-rank; base=h2

\section{第四十二章——论宗教的智慧：基督的名，除祂自己和圣父以外，无人知晓。}

我现在要说明什么是智慧的宗教，或宗教的智慧。天主在起初，在创造世界之前，从祂永恒的本源，并从神圣而永生的圣神，为自己生了一位圣子：不败坏的、忠信的、与圣父的卓越和尊威相称的。祂是德能，祂是理性，祂是天主的圣言，祂是智慧。这位工匠（如赫尔墨斯所说）与谋士（如西比拉所说）一同规划了这个世界卓越而奇妙的构造。总之，在众天使中（这些天使是由同一天主从祂的呼吸中形成的），唯独祂被接纳分享祂至高权能，唯独祂被称为天主。因为万物都是借着祂，而没有祂就没有任何东西。最后，柏拉图不完全作为哲学家，而是作为先知，谈论了第一位天主和第二位天主，也许在此追随了特里斯墨吉斯图斯；我将他的话从希腊文译出并附在下面：\Quote{万有的主与创造者，我们以为应称为天主的那一位，创造了一位第二位天主，是可见可感的。但我说的可感，并不是说祂自己接受感觉，而是说祂引起感觉与视觉。因此，当祂造了这位第一位、唯一的一位、独一无二的一位之后，祂在祂眼中看为极其卓越，充满了一切美善。}西比拉也说，万有的引导者是天主所造的，另一位又说，\par

\Quote{必须认识天主子——天主，}\par

正如我在我的著作中所列举的那些例子所宣告的那样。众先知，充满了圣神的灵感，预言了祂；其中尤其是撒罗满在《智慧书》中，以及他的父亲——神圣颂歌的作者——这两位最著名的君王，比特洛伊战争早一百八十年，他们作证说祂是天主所生的。祂的名，除了祂自己和圣父以外，无人知晓，正如若望在《默示录》中所教导的。\BibleRef{默 19:12}赫尔墨斯说，祂的名是凡人的口所不能言说的。然而，人们用两个名字称呼祂：耶稣，即救主；基督，即君王。祂被称为救主，因为祂是通过祂而信从天主之人的健康与救恩。祂被称为基督，因为在这时期结束时，祂将亲自从天降来审判世界，并在复活死者后，为自己建立永恒的王国。\par

% structure: p0109 source=h2 target=section reason=relative-heading-rank; base=h2

\section{第四十三章——论耶稣基督的名，及祂的双重诞生。}

但是，恐怕你心中对为何称祂为耶稣基督有任何疑问，祂是在世界之前由天主所生，又在三百年前由人所生，我将向你简要说明理由。同一位是天主之子，也是人之子。因为祂两次诞生：第一次由天主在精神上诞生，在世界起源之前；后来在\Person{奥古斯都}统治时期，由人在肉身中诞生；与此事相关的是一个显赫而伟大的奥迹，其中包含了人类的救恩、至高天主的宗教，以及一切真理。因为当那可诅咒且不虔敬的敬拜诸神的行为首先藉着魔鬼的诡计潜入时，天主的宗教只保留在希伯来人中，他们不是依据任何法律，而是依照祖先的方式，遵守历代相传的敬拜，直到他们由\Person{梅瑟}（众先知中第一位）率领出离\Place{埃及}的时候；藉着\Person{梅瑟}，法律从天主赐给了他们；他们后来被称为犹太人。因此，他们服侍天主，被法律的锁链束缚。但他们也逐渐偏离到亵渎的仪式中，开始了对异邦神祇的敬拜，离弃了祖先的敬拜，向无知的偶像献祭。因此，天主派遣充满圣神的先知到他们那里，以他们的罪过责备他们，宣讲悔改，以将临的惩罚威胁他们，并宣布若他们坚持同样的过错，祂将派遣另一位作为新法律的携带者；并要将这忘恩负义的百姓从他们的产业中移除，从外邦中为自己召集一个更忠信的百姓。但他们不仅坚持己路，甚至杀害了这些使者。因此，天主因这些行为而谴责他们；祂不再向顽梗的百姓派遣使者，而是派遣了自己的儿子，呼召万民归向天主的恩宠。然而，祂并没有将他们（尽管他们不虔敬和忘恩负义）排除在救恩的希望之外：祂首先派遣祂到他们那里，以便如果他们偶然服从，就不致失去所领受的；但如果他们拒绝接受他们的天主，那么，继承人被移除后，外邦人将承受产业。因此，至高圣父命令祂降临到地上，穿上人的身体，使祂藉着肉身所受的苦难，不仅以言语，也以行动教导德行和忍耐。因此，祂第二次作为人由童贞女所生，没有父亲，好使祂第一次属灵诞生（只由天主所生，成为神圣的精神）与第二次肉身的诞生（只由母亲所生，成为圣洁的肉身）相称，以便藉着祂，那已屈服于罪恶的肉身得以脱离毁灭。\par

% structure: p0111 source=h2 target=section reason=relative-heading-rank; base=h2

\section{第四十四章．从先知证明基督的双重诞生。}

这些事应这样发生，正如我所陈述的，先知们早已预言了。在\Person{撒罗满}的著作中这样写着：\Quote{一位童贞女的子宫被坚固，就怀了孕；一位童贞女受了孕，在极大的怜悯中成为母亲。}在\Person{依撒意亚}书中（\BibleRef{依 7:14}）这样写着：\BibleQuote{看，一位童贞女要怀孕生子，你要给祂起名叫\OriginalTerm{厄玛奴耳}{Immanuel}；}这名字翻译出来，就是\Quote{天主与我们同在}。（\BibleRef{玛 1:23}）因为祂在世上与我们同在，当祂取了肉身时；祂在天主内同样是天主，在人间同样是人。祂既是天主又是人，先知们早已宣告了。祂是天主，\Person{依撒意亚}（\BibleRef{依 45:14}）这样宣告：\BibleQuote{他们要向你俯伏，向你恳求；因为天主在你内，我们却不知道，祂是以色列的天主。凡反对你的人都要羞愧蒙羞，归于混乱。}还有\Person{耶肋米亚}（\BibleRef{巴 3:35}）：\BibleQuote{这是我们的天主，无人可与祂相比；祂发现了知识的全部道路，并将它赐给了祂的仆人\Person{雅各伯}，和祂所爱的\Place{以色列}。此后祂在地上显现，住在人间。}同样，祂是人，同一\Person{耶肋米亚}说：\Quote{祂是人，有谁认识祂？}\Person{依撒意亚}也这样说（\BibleRef{依 19:20}）：\BibleQuote{上主要给他们派遣一个人，拯救他们，并以公义医治他们。}\Person{梅瑟}本人也在\WorkTitle{户籍纪}中：\BibleQuote{由\Person{雅各伯}要出现一颗星，由\Place{以色列}要兴起一人。}因此，为这缘故，祂既是天主，就取了肉身，成为天主与人之间的中保，战胜了死亡，藉着祂的引导将人带领到天主面前。\par

% structure: p0113 source=h2 target=section reason=relative-heading-rank; base=h2

\section{第四十五章．从圣经证明基督的能力和作为。}

我们已经说了祂的诞生；现在让我们谈谈祂的能力和作为。当祂在人间施行这些时，犹太人看到它们伟大而奇妙，便以为它们是藉着魔术的影响力成就的，却不知道祂所行的一切都已被先知预言了。祂赐力量给病人，以及被各种疾病缠身的人，不是藉着任何治疗的药物，而是瞬间藉着祂言语的力量和大能；祂使软弱的人复原，使瘸子行走，使瞎子看见，使哑巴说话，使聋子听见；祂洁净了被污染和不洁的人，使那些被魔鬼攻击而疯狂的人恢复清醒，使那些已死或被埋葬的人重获生命和光明。祂还以五个饼和两条鱼喂饱并满足了五千人。祂也在海上行走。祂还在风暴中命令风平静，立刻就有了风平浪静；所有这一切，我们发现在先知书和西比尔女巫的诗歌中都有预言。\par

当因这些奇迹而有一大群人归向祂，并正如祂确实所是，相信祂是天主之子、由天主所派遣时，犹太人的司祭和首领满怀嫉妒，同时又因祂责备他们的罪恶和不义而被激怒，便合谋杀害祂。这件事在约一千多年前，撒罗满在\WorkTitle{智慧篇}中已经预言，用了这些话：\Quote{让我们欺骗义人，因为他对我们不合意，并责备我们违犯法律。他自夸拥有天主的智识，并称自己为天主的儿子。他被造来指责我们的思念：我们看他一眼都感到痛苦；因为他的生活与别人不同，他的道路是另一种方式。我们被他视为轻浮之人；他如同远离污秽一般远离我们的道路；他极力称赞义人的结局，并夸耀天主是他的父亲。让我们看看，他的话是否真实；让我们证明他会有何种结局；让我们用斥责和折磨来考验他，好知道他的温良，并证明他的忍耐；让我们判他受可耻的死亡。他们这样想象，便走差了路；因为自己的愚昧使他们眼瞎，他们不明白天主的奥秘。}\par

因此，他们对自己所读的这些经文不加留意，竟煽动百姓如同对付一个不敬神的人，以致他们捉拿祂并带去受审，并以亵渎的言词要求处死祂。但他们指控祂的罪状正是这一点：祂说祂是天主之子，并因在安息日治病而违背了法律——而祂说自己并没有违背，反倒是成全了。当那时以总督身份管辖叙利亚的般雀·比拉多看出这案件不属于罗马审判官的职责范围时，便将祂送到黑落德分封侯那里，并允许犹太人自己按照他们的法律审判；这些人获得惩罚祂罪行的权力后，便判处祂钉十字架，但先是鞭打祂，用手击打祂，给祂戴上荆棘冠冕，向祂脸上吐唾沫，给祂吃醋和苦胆；在这些事中，没有一句话从祂口中发出。然后，行刑者拈阄分祂的袍子和外衣，将祂悬在十字架上，钉住祂，尽管第二天他们就要庆祝逾越节，即他们的节日。这罪行随即被异象所跟随，好叫他们明白自己所行的不敬；因为在祂断气的同一时刻，发生了大地震，太阳隐没，白昼变成黑夜。\par

% structure: p0117 source=h2 target=section reason=relative-heading-rank; base=h2

\section{第四十六章——从先知证明基督的受难与死亡已被预言。}

众先知已预言这一切将会如此发生。依撒意亚这样说：\BibleRef{依 50:5}\Quote{我并非悖逆，也不抵抗；我把我的背转给鞭打者，我的面颊转给掌击者；我没有掩面躲避唾污。}同一先知论及祂的沉默说：\BibleRef{依 53:7}\Quote{我被牵去如同待宰的羊，又像羔羊在剪毛者前缄默，祂也没有开口。}达味也在第三十四篇圣咏中说：\Quote{卑贱之徒聚集攻击我，他们不认识我；他们分散却不懊悔；他们试探我，向我切齿。}同一圣咏家在第六十八篇中论及食物和饮料说：\Quote{他们拿苦胆给我作食物，我渴时，他们拿醋给我喝。}关于基督的十字架又说：\Quote{他们穿透了我的手和我的脚，数算了我的所有骨头；他们瞪着眼看我，分了我的衣服，为我的长衣拈阄。}梅瑟也在《申命纪》中说：\BibleRef{申 28:66}\Quote{你的性命必悬于眼前，你昼夜恐惧，不得保全性命。}又在《户籍纪》中说：\BibleRef{户 23:19}\Quote{天主不像人那样犹豫，也不像人子那样受威胁。}匝加利亚也说：\BibleRef{匝 12:10}\Quote{他们要仰望我所刺透的那位。}亚毛斯论及日蚀这样说：\BibleRef{亚 8:9}\Quote{在那一天，——上主说——太阳要在正午下落，晴朗的白昼变为黑暗；我要使你们的节庆变为悲哀，你们的歌曲变为哀歌。}耶肋米亚论及祂受苦的耶路撒冷城也说：\BibleRef{耶 15:9}\Quote{她的太阳在白天落下，她受羞辱、被咒骂，其余的人我将交给刀剑。}这些话并非空言。不久之后，韦斯帕芗皇帝征服了犹太人，用刀剑和火荒废了他们的土地，围困他们，使他们饥荒，推翻了耶路撒冷，将俘虏示众游街，并禁止其余人返回故土。这些事是天主因基督被钉十字架而行的，正如祂先前在圣经中对撒罗满所宣告的：\BibleRef{列上 9:7}\Quote{以色列必成为万民中的毁灭和笑柄，这殿必成为废墟；凡经过的人都要惊愕，说：上主为什么向这地和这殿行了这样的恶事？他们必说：因为他们离弃了上主他们的天主，迫害了他们那位为天主所深爱的君王，并以极大的羞辱将祂钉在十字架上，因此上主使这一切灾祸临到他们。}因为杀害那为他们的救恩而来的主，还有什么不配得的呢？\par

% structure: p0119 source=h2 target=section reason=relative-heading-rank; base=h2

\section{第四十七章——论耶稣基督的复活、宗徒的派遣及救主升天。}

此后，他们将祂的身体从十字架上取下，葬在坟墓里。但在第三日黎明前，发生了地震，堵住坟墓的石头被挪开，祂复活了。坟墓中除了包裹祂身体的殓布之外，什么也没有。先知们早已预言祂将在第三日复活。\Person{达味}在\BibleBook{}{圣咏}第十五篇中说：\BibleQuote{祢不会将我的灵魂遗弃在地狱，也不让祢的圣者见到腐朽。} \Person{欧瑟亚}也说：\BibleQuote{这是我的智慧之子，因此他必不长久停留在众子的痛苦中；我要从坟墓的手中赎他。死亡啊，你的审判在哪里？你的毒刺在哪里？} 同一位先知又说：\BibleRef{欧 6:2} \BibleQuote{两天后祂必使我们复生，第三天祂必使我们兴起。}\par

因此，祂复活后去了\Place{加里肋亚}，再次召集了因恐惧而逃散的众门徒；向他们颁发了祂要遵守的命令，安排了向普世宣讲福音的事宜，然后向他们嘘气，赐予圣神，并赋予他们行奇迹的能力，使他们能以言以行造福人类。最后，在第四十天，祂被云彩接去，升到圣父那里。\Person{达尼尔}先知早已预示此事：\BibleRef{达 7:13} \BibleQuote{我在夜间的异象中观望，看见一位相似人子者，乘着天上的云彩而来，走向万古常存者，侍奉祂的人将祂引到祂面前。于是赐给祂统治权、尊荣和国度，各民族、各邦国、各语言的人都侍奉祂；祂的王权是永远的，永不消逝，祂的国度永不灭亡。} \Person{达味}在\BibleBook{}{圣咏}第一〇九篇中也说：\BibleQuote{上主对我主说：你坐在我右边，等我使你的仇敌做你的脚凳。}\par

% structure: p0122 source=h2 target=section reason=relative-heading-rank; base=h2

\section{第四十八章——论犹太人的被剥夺继承权与外邦人的被接纳}

既然祂坐在天主的右边，即将践踏那些折磨祂的仇敌，当祂来审判世界时，显然犹太人已无希望，除非他们回头悔改，并洗净他们染上的血污，开始寄望于他们所否认的那一位。因此，\Person{厄斯德拉}这样说：\BibleQuote{这逾越节是我们的救主和避难所。你们要思量，让它深入你们的心：我们必须以预像贬抑祂；此后我们寄望于祂。}\par

犹太人因拒绝基督而被剥夺继承权，而我们这些外邦人被接纳取代他们，这由圣经所证明。\Person{耶肋米亚}这样说：\BibleRef{耶 12:7} \BibleQuote{我离开了我的家，把我的产业交在敌人的手中。我的产业对我竟像林中的狮子，它向我发出声音，因此我憎恨它。} \Person{马拉基亚}也说：\BibleRef{拉 1:10} \BibleQuote{我不喜欢你们，——万军的上主说——也不接受你们手中的祭品。因为从日出之地到日落之处，我的名在异民中大受尊崇。} \Person{依撒意亚}也这样说：\BibleRef{依 66:18} \BibleQuote{我要聚集万民和万国，他们要来看我的荣耀。} 同一位先知在另一处，以圣父对圣子的口吻说：\BibleRef{依 42:6} \BibleQuote{我，上主，因正义召叫了你，我必提携你的手，保护你，立你作人民的盟约，外邦人的光明，为开启盲人的眼目，从牢狱中领出被囚的人，从暗牢中领出坐在黑暗中的人。}\par

% structure: p0125 source=h2 target=section reason=relative-heading-rank; base=h2

\section{第四十九章——论天主是唯一的}

如果犹太人被天主弃绝，正如神圣典籍所启示的信仰所示，而外邦人如我们所见的被接纳，并摆脱了今世的黑暗和魔鬼的锁链，那么，除非人追随真正的宗教和真正的智慧——那就是在基督内——否则别无希望；不认识祂的人永远远离真理和天主。犹太人或哲学家们不要以至尊天主自夸。凡不认识圣子的，不能认识圣父。这就是智慧，这就是至尊天主的奥秘。天主愿意人们藉着圣子认识并敬拜祂。为此，祂预先派遣先知们宣布圣子的来临，以便当预言应验在祂身上时，人们能相信祂既是天主子，也是天主。\par

然而，不可持有两个天主的见解，因为圣父与圣子是一体。既然圣父爱圣子并将一切交给祂，而圣子忠诚地服从圣父，除了圣父所做的，别无所愿，那么显而易见，如此亲密的关系不可分离，以至于不能说他们为二，在他们内只有一个实体、一个旨意和一个信德。因此，圣子藉圣父而存在，圣父藉圣子而存在。两者应受同等的尊荣，犹如一位天主，并且在敬拜两者时应有所区分，使得这区分本身因不可分割的合一之纽带而结合。谁若将圣父与圣子分离，或将圣子与圣父分离，就是弃绝自己的一切。\par

% structure: p0128 source=h2 target=section reason=relative-heading-rank; base=h2

\section{第五十章——论天主为何取可朽的身体并承受死亡}

还须回答那些认为天主穿上可朽的身体、服从于人、忍受侮辱、甚至遭受酷刑和死亡是不得体和不合理的人。我要说出我的见解，并且将尽我所能，用几句话概括一个巨大的主题。我认为，凡教导者，他自己应当实践他所教导的，以促使人们服从。因为若他不实践，便会减损对其训诫的信德。因此需要榜样，使所给的训诫具有坚固性；若有人倔强不服，声称这些训诫无法在实行中完成，教导者便可以用实际行动反驳他。因此，一套教导体系当它仅以言语传授时不能是完美的；只有当它由行为完成时才成为完美。\par

既然基督被派遣到人间作为德性之师，为了祂教导的完善，祂理当既行事又教导。但如果祂没有取人身，祂就不能实践祂所教导的——即不动怒、不贪财富、不燃起情欲、不畏惧痛苦、蔑视死亡。这些显然是德行，但没有肉身就无法做到。因此祂取了肉身，正是为了：既然祂教导必须克服肉身的欲望，祂亲自率先实践，使无人能以肉身的软弱为借口。\par

% structure: p0131 source=h2 target=section reason=relative-heading-rank; base=h2

\section{第五十一章 论基督在十字架上的死亡}

现在我要谈论十字架的奥迹，以免有人碰巧说：如果祂必须死，也当是一种明显不名誉、不光彩的死，而非某种有体面的死。我确实知道，许多人厌恶十字架之名，便回避真理，尽管其中包含着极大的合理性和力量。因为祂被派遣正是为此目的：为最卑微的人打开救恩之路，祂使自己谦卑，为要释放他们。因此祂承受了那种通常加于卑微者身上的死亡，使众人都有效法的机会。此外，因为祂将要复活，祂的身体不容有任何部位被残割或骨头折断——这与被斩首之人的遭遇相反。因此十字架被选中，它保存身体与骨头完好无损，以供应复活。\par

此外还加上这些理由：祂既决意受苦和受死，理当被高举。因此十字架既在事实上也在象征上高举了祂，使祂的尊威和能力连同苦难一同为人所知。因为当祂在十字架上伸展双手时，祂显然向东方和西方张开祂的翅膀，天下的万民从世界的两端都可以聚集并安息在祂之下。这标记何等重大、有何等权能，是显而易见的，因为所有魔鬼的军旅都因这标记而被驱逐和溃散。正如祂自己在受难前以言语和命令驱魔，如今，借着同一苦难的名号和标记，附在人身上的不洁之神，在备受折磨之中被赶出，并承认自己是魔鬼，顺从于折磨他们的天主。那么，希腊人看到他们的神明——他们也承认这些是魔鬼——被人借着十字架制服时，还能指望从他们的迷信和智慧中得到什么呢？\par

% structure: p0134 source=h2 target=section reason=relative-heading-rank; base=h2

\section{第五十二章 人的救恩之望在于认识真天主，以及异教徒对基督徒的仇恨}

因此，人只有一重生命之望、一个安全港口、一个自由避难所：就是放下纠缠他们的错谬，睁开他们心灵的眼睛，认识天主——独有祂是真理的居所；轻视地上的事物，视泥土所造之物为无，视在天主眼中是愚拙的哲学为无；并接受真智慧即宗教，成为永生的继承者。然而，他们与其说反对真理，不如说反对自身的得救；一旦听到这些事，便如同厌恶某种不可赎的邪恶一样憎恶它们。他们甚至不肯听：他们以为自己的耳朵被不虔敬玷污了；现在他们不仅不停止辱骂，还用最侮辱的言辞攻击；而且，若得了权力，就像对待公敌一样迫害他们——甚至比对待敌人更甚；因为敌人一旦被征服，便以死亡或奴役来惩罚；放下武器之后便不再施加酷刑，尽管那些人理当遭受一切——他们企图使虔敬在刀剑之间仍有立足之地。\par

残忍与清白相结合是闻所未闻的，也不配得上胜利的敌人的身份。如此愤怒的巨大原因是什么？无疑，因为他们不能在理性的基础上争论，便以暴力推进自己的主张；在不了解对象的情况下，他们谴责那些拒绝就自身清白事实进行辩护的人是极其有害的。而且，他们认为仅让那些他们无理憎恨的人迅速简单地死去还不够，还用精巧的酷刑来折磨他们，以满足他们的仇恨——这仇恨并非由任何过错产生，而是由真理产生，真理为那些生活邪恶的人所憎恨，因为他们恼火存在一些不为他们的行径所悦的人。他们想方设法毁灭这些人，为能在没有任何见证人的情况下毫无约束地犯罪。\par

% structure: p0137 source=h2 target=section reason=relative-heading-rank; base=h2

\section{第五十三章 查考并驳斥仇恨基督徒的理由}

但他们说做这些是为了保卫他们的神。首先，如果他们是神，且拥有任何能力和影响，他们就不需要人的保卫和保护，而是会明显地自卫。或者，如果他们连自己的伤害都不能报复，人又如何能期望从他们那里得到帮助呢？因此，想要成为神的复仇者是徒劳而愚蠢的，只是由此更显明他们的不信。因为承担他所崇拜之神的保护者，承认了那神的无能；但如果他因为认为那神有能而崇拜他，他就本不应想保护他，而应被他保护。因此我们做得对。因为当那些假神的保卫者，即反抗真天主的人，在我们身上迫害祂的名时，我们既不反抗行为也不反抗言语，而是以温良、缄默和忍耐忍受任何残酷所能对我们施行的。因为我们信赖天主，期望日后将临的报应。这信赖并非没有根据，因我们有时听说、有时也亲眼看见所有敢于犯此罪者的悲惨结局。没有人能不受惩罚地侮辱天主；凡不愿借言语学习的人，便借自己的惩罚知道了谁是真天主。\par

我想知道，当他们强迫人违背意愿祭献时，他们自己心里有什么理由，或者他们向谁祭献。如果是献给神，那么由不受神喜悦的人所做、因伤害而勒索、因痛苦而强加的，就不是崇拜，也不是可接受的祭献。但如果是做给那些被强迫的人，那显然不是什么恩惠，因为没人会接受它，他甚至宁愿死。如果你召唤我去行善，为何用恶来邀请我？为何用鞭打而不用言语？为何不用论证而用身体折磨？由此显然，你所引诱我去的，不是自愿而是拖着我拒绝的那个，是恶。违背他人意愿而想为他谋福，是何等愚蠢！如果有人迫于苦难试图求死，你若从他手中夺下剑、割断绞索、从悬崖边拉回他、或倒掉毒药，你能自夸为那人的救命恩人吗？而你认为救了他的人，并不感谢你，反而认为你对他做了坏事，阻止了他所渴望的死，不让他抵达终点和从劳苦中安息。因为恩惠不应按行为性质衡量，而应按接受者的感受。你为何要把对我有害之事算作恩惠？你希望我崇拜你的神，而我视之为对我有害。如果那是善，我不嫉妒。你独自享受你的善吧。你没有理由想救助我的错误，这错误是我凭我的判断和意愿承担的。如果是恶，你为何拖我去分享恶？你自己用你的命运吧。我宁愿在行善中死去，也不愿在恶中活着。\par

% structure: p0140 source=h2 target=section reason=relative-heading-rank; base=h2

\section{第五十四章——论崇拜天主的宗教自由。}

这些话也许说得公正，但有谁会听呢？当那些狂暴不羁的人认为，如果人事中有任何自由，他们的权威就会受到削弱。然而唯独宗教是自由居处之地。因为它本是最自愿的事，不能对任何人施加必要性，以迫使他崇拜他所不愿崇拜的。也许有人会假装，但他不能愿意。总之，有些人因恐惧折磨或经受酷刑而同意了可憎的祭献：他们从不会自愿做那些出于不得已之事；但当再次有机会、自由恢复时，他们便再次归向天主，以祈祷和眼泪向祂赎罪，后悔的不是意志（他们本没有），而是所受的不得已；对做补赎的人，宽恕不会被拒绝。那么，玷污身体的人又成就了什么？因为他不能改变意志。\par

但事实上，心智薄弱的人，如果他们将任何有勇气的人引诱去祭献给他们，便以难以置信的兴奋狂妄得意，好像使敌人臣服于轭下。但若有人既不因威胁也不因酷刑而畏惧，选择宁可为信仰而舍弃生命，残忍就使出全部心机对付他，策划可怕而难以忍受的事；因为他们知道为天主的事业而死是光荣的，并且我方是胜利，如果我们在战胜施刑者后为信仰和宗教献出生命，他们也就自己努力要战胜我们。他们不处死我们，而是寻找新的、前所未闻的折磨，使肉体的脆弱屈服于痛苦；若不屈服，他们就延后进一步的惩罚，并对伤口加以悉心照料，使伤疤尚新时重复的酷刑带来更多疼痛；当他们以无辜者施行这种酷刑时，显然认为自己虔诚、公义、虔诚（因为他们喜欢这样的祭献给他们的神），却称另一些人为不虔诚和绝望。这是何等颠倒：受罚者虽无辜，却被称为绝望和不虔诚，而施刑者反被称为公义和虔诚！\par

% structure: p0143 source=h2 target=section reason=relative-heading-rank; base=h2

\section{第五十五章——外教人指责正义之人在跟随天主上的不虔诚。}

但有人说，那些不喜欢从祖先传下来的公共宗教礼仪的人，应受公正且应得的惩罚。假如那些祖先在从事虚妄的宗教礼仪时是愚昧的，正如我们先前所证明的，难道我们就应被禁止遵循更真、更好的事物吗？为什么我们剥夺自己的自由，并像被捆绑一样，成为他人错误的奴隶？请允许我们有智慧，请允许我们探求真理。然而，如果他们乐意捍卫他们祖先的愚昧，为什么埃及人得以逃脱，他们崇拜牲畜和各种野兽为神明？为什么诸神自己成为喜剧表演的对象？为什么那以最机智的方式嘲笑他们的人反而受尊敬？为什么哲学家们被聆听，他们要么说没有神，要么说即使有神，他们对人间事务也不关心、不照管，或者论证根本没有统治世界的天主眷顾？\par

但在所有人中，唯独那些跟随天主和真理的人被判定为不虔敬。既然这就是正义与智慧，他们却指控它为不虔敬或愚昧，不明白是什么欺骗了他们，当他们称恶为善、称善为恶。许多哲学家，尤其是\Person{柏拉图}与\Person{亚里士多德}，确实谈论了许多关于正义的事，声称并极度赞扬这一德行，因为它给予每人应得的，因为它在一切事上维持公平；而其他德行仿佛是静默的、封闭于内的，唯独正义既不只关心自身，也不隐藏，而是完全展现于外，并乐于施惠，以帮助尽可能多的人——仿佛正义本当只存在于法官和那些身居任何权位者中，而不是在所有人中。\par

然而，没有一个人，甚至最卑微者和乞丐，不能行正义。但因他们不知道正义是什么、从何而来、其运作方式如何，他们便将这至高德行（即众人的共同善）只归于少数人，并说它不追求任何自身的利益，而只追求他人的利益。\Person{嘉尼迪斯}，一位极具才华和洞察力的人，起来反驳他们的言论并推翻那无坚实基础的正义，并非没有理由；这并非因为他认为正义应受指责，而是为表明正义的捍卫者并未提出关于正义的任何可靠或确定的论证。\par

% structure: p0147 source=h2 target=section reason=relative-heading-rank; base=h2

\section{第五十六章——论正义，即对真天主的敬拜。}

因为如果正义是对真天主的敬拜（论公平，有什么比承认天主为父、尊崇祂为主、服从祂的法律或诫命更公平？论尊荣，有什么比这更虔敬？论安全，有什么比这更必要？），那么哲学家们并不认识正义，因为他们既不承认天主本身，也不遵守祂的敬拜和法律；正因如此，他们本可被\Person{嘉尼迪斯}所驳斥，他的辩论大意是：没有自然的正义，因此所有动物都凭本性引导维护自身利益，所以正义若促进他人利益而忽略自身利益，就应称为愚昧。但如果所有掌权者，尤其是统治全世界的罗马人，愿意遵循正义，将他们用武力和武器夺取的财物归还各人，他们将回到茅舍和贫困的境地。如果他们这样做，他们固然会是正义的，但必然被视为愚昧，因为他们为了他人的利益而损害自己。再者，如果有人因错误将金子当作黄铜、银子当作铅来卖，而某人因需要不得不购买它，他会隐瞒知识并以低价购买，还是会告知卖者其价值？如果他告知，他将明显被称为正义，但也是愚昧的，因为施益于人而损害自己。但在伤害案中容易判断。如果他面临生命危险，以致必须要么杀他人要么死，他将怎么做？可能的情况是，遭遇船难后，他发现一个虚弱者抱着一块木板；或军队战败逃亡时，他发现一个骑马的伤者：他是否会将那人从木板上推下、从马上推下，以便自己能逃脱？如果他愿意正义，他不会这样做；但他也将被视为愚昧，因保全他人性命而丧失自己的性命。如果他这样做，他确实会显得有智慧，因为他会维护自身利益，但他也是邪恶的，因为他会行不义。\par

% structure: p0149 source=h2 target=section reason=relative-heading-rank; base=h2

\section{第五十七章——论智慧与愚昧。}

这些言论确实犀利，但我们很容易回应它们。因为对名称的模仿导致如此显现。正义与愚蠢相似，然而它并非愚蠢；同时，恶意与智慧相似，然而它并非智慧。但正如那种恶意在维护自身利益时是聪明而机敏的，它并非智慧，而是狡猾与狡诈；同样，正义也不应被称为愚蠢，而应称为纯真，因为义人必须是智慧的，而愚人则是不义的。因为无论理性还是自然本身都不允许义人不智慧，因为显然义人只做正确和良善的事，并且总是避免乖僻和邪恶的事。然而，除了那智慧的人，谁能区分善与恶、邪僻与正直呢？但愚人行恶，因为他不知善恶。因此他做错事，因为他无法区分邪僻与正直。所以，正义不适合愚人，智慧也不适合不义之人。因此，那个没有将溺水者从木板上推开、没有将受伤者从马上推开的人并不是愚人，因为他避免了伤害——伤害是一种罪；而避免罪是智者的本分。\par

但他之所以初看显得愚蠢，是因为他们以为灵魂与身体一同消灭；为此他们将一切利益归于今生。因为如果死后没有存在，那么显然，牺牲自己的利益去顾惜他人的生命，或更多顾及他人的利益而非自己的利益，是愚蠢的行为。如果死亡毁灭了灵魂，我们就必须努力活得更长久，更有利于自己；但如果死后有永生和福乐的生命，那么正义而智慧的人必定会轻视这肉体的存在及一切尘世财物，因为他知道自己将从天主那里得到何种赏报。因此，让我们持守纯真，持守正义，甘愿承受愚蠢的外表，好能持守真正的智慧。如果世人认为宁愿受酷刑和死亡而不向诸神献祭并安然脱身是愚昧和愚蠢的，但让我们务必以一切德性和一切忍耐努力显明对天主的忠诚。不要让死亡惊吓我们，不要让痛苦征服我们，以致我们心灵的力量和坚定不能保持不动摇。让他们称我们为愚人吧，但他们自己却是最愚蠢的，瞎眼、迟钝、如同羊群；他们不明白离开永生的天主而俯伏敬拜尘世之物是致命的事；他们不知道永恒的惩罚等待着那些敬拜无灵之像的人；而那些为敬拜和尊崇真天主而不拒绝酷刑和死亡的人，将获得永生。这就是至高信德；这就是真智慧；这就是完美的正义。愚人如何判断、轻浮之人如何想法，与我们无关。我们应等待天主的审判，以便日后审判那些审判过我们的人。\par

% structure: p0152 source=h2 target=section reason=relative-heading-rank; base=h2

\section{论对天主的真实敬拜与祭献。}

我已经论及正义，其性质如何。接下来我要表明什么是献给天主的真实祭献，什么是最公义的敬拜祂的方式，以免有人认为牺牲、乳香或贵重礼物是天主所愿的。祂既不受饥饿、干渴、寒冷及一切尘世欲望的支配，因此也不会使用所有呈奉在庙宇和向地上诸神献上的这些东西；正如有形的祭品对于有形之物是必要的，同样，无形的祭品对于无形之体显然是必要的。但天主并不需要那些祂赐予人使用的东西，因为全地都在祂权下；祂不需要圣殿，因为世界是祂的住所；祂不需要偶像，因为祂眼目和心智都无法理解；祂不需要地上的光，因为祂能点燃太阳的光以及其它星辰供人使用。那么天主向人要求什么呢？不过是纯净而神圣的心灵敬拜。因为人手所造的或人以外的那些东西，是无灵的、脆弱的、不讨人喜欢的。这才是真实的祭献，不是从胸中而是从心中产生的；不是用手而是用心灵奉献的。这才是悦纳的祭品，是心灵献上自己的祭品。因为祭品有何用？乳香有何用？衣袍有何用？银子有何用？金子有何用？宝石有何用？——如果敬拜者没有纯洁的心灵。因此，天主只要求正义。在此有祭献；在此有天主的敬拜，关于这点我现在必须论述，并表明正义必然包含在哪些行为中。\par

% structure: p0154 source=h2 target=section reason=relative-heading-rank; base=h2

\section{论生活方式与世界之初。}

哲学家和诗人都知道人的生命有两条道路，但两者以不同的方式提出。哲学家们希望一条是勤劳之路，另一条是懒惰之路；但他们在这方面说得不够准确，因为他们只将这两条路指向此生的利益。诗人说得更好，他们说一条是义人的道路，另一条是不义之人的道路；但他们在这事上犯错，因为他们说这两条路不在今生，而在下界阴间。我们所说的则显然更正确，我们说一条是生命之路，另一条是死亡之路。然而我们在此说，有两条路；但右边的那条路，义人在其上行走，并不通往\Place{埃律西昂}，而是通往天堂，因为他们成为不朽的；左边的那条路通往\Place{塔尔塔罗斯}，因为不义之人被判受\OriginalTerm{永罚}{eternal tortures}。因此，我们应当持守那通向生命的正义之路。如今，正义的首要义务是承认天主为父，敬畏祂为主，爱祂如父。因为那生了我们、以生命的灵气赋予我们活力、滋养并保存我们的同一位，对我们不仅有如父亲，也有如主人，有纠正我们的权威及生杀大权；因此人当向祂献上双重的尊敬，即爱与敬畏相结合。正义的第二项义务是承认人人为兄弟。因为如果同一天主创造了我们，并使所有人在平等的条件下趋向正义和永生，那么显然我们因兄弟情谊而联合；不承认这一点的人是不正义的。但这一邪恶的根源——它撕裂了人与人之间的相互交往，撕裂了亲属关系的纽带——源于对真天主的无知。因为那不知道恩惠之源的人，绝不可能成为善人。因此，从那时起，当众多神祇被祝圣并受人敬拜时，正义——正如诗人所叙述的——被驱逐，一切盟约被摧毁，人类正义的团契被摧毁。于是，每个人都只为自己打算，以为强权即公理，伤害别人，以诡计攻击，以奸诈欺骗，以他人的不便来增加自己的利益，不宽容亲戚、子女或父母，为毁灭人而调制毒杯，以刀剑拦路，骚扰海上，放纵情欲，无论情欲引领到哪里——简言之，认为没有任何东西是神圣的，不为他那可怕的欲望所侵犯。当这些事情发生时，人们便为自己设立法律，以促进公共利益，同时保护自己免受伤害。但法律的恐惧并未压制罪行，而是抑制了放荡。因为法律能惩罚罪行，却不能惩罚\OriginalTerm{良心}{conscience}。因此，以前公开做的事开始秘密地进行。正义也因诡计而被规避，因为那些自己执掌法律执行的人，被礼物和赏报所腐化，拿判决做交易，或让恶人逃脱，或毁灭善人。在这些之上，又添上了纷争、战争和互相掠夺；法律被践踏，暴力行事的权力毫无约束地被攫取。\par

% structure: p0156 source=h2 target=section reason=relative-heading-rank; base=h2

\section{第60章——论正义的义务}

当人类的事务处于这种状况时，天主怜悯了我们，向我们启示并显示了自己，好叫我们能在祂内学习宗教、信德、纯洁和仁慈；好叫我们放下先前生活的错误，与天主自身一起认识我们自己——我们曾因不虔诚而与祂分离——并且我们可以选择那将人类事务与天上事务结合起来的\OriginalTerm{神圣法律}{divine law}，是主亲自将它交付给我们；藉此法律，我们被缠绕的一切错误，连同虚妄而不敬的迷信，都能被除去。因此，我们对人应尽的本分，正是由那同一神圣法律所规定的，它教导：你施予人的一切，就是施予天主。但\OriginalTerm{正义的根源}{root of justice}，以及公平的全部基础，就在于：你不应做你不愿承受的事，而应以你自己的感受来衡量他人的感受。如果承受伤害是不愉快的，而做伤害的人似乎是不义的，那么把你对自己所感受的转移到他人的身上，把你对他人所判断的转移到自己的身上，你就会明白，如果你伤害别人，就像别人伤害你一样，你的行为是不义的。如果我们思考这些事，我们就能保持\OriginalTerm{无辜}{innocence}，可以说，正义的第一步就包含在其中。因为首先是不伤害；其次是服务他人。正如在未开垦的土地上，在你开始播种之前，必须先清除田地，拔除荆棘，砍掉树根的根须；同样，必须先驱逐我们灵魂中的恶习，然后再植入\OriginalTerm{美德}{virtues}，如此，由天主的圣言所产生的永生之果才能生长出来。\par

% structure: p0158 source=h2 target=section reason=relative-heading-rank; base=h2

\section{第61章——论情欲}

有三种\OriginalTerm{情欲}{passions}，或者可谓三种仇神，它们在人的灵魂中激起如此巨大的骚动，有时迫使人们以这样的方式犯罪，以致他们既不顾及自己的名誉，也不顾个人安全：这就是愤怒（渴求报复）、贪财（渴望财富）、情欲（追求快乐）。我们必须首先抵制这些恶习；必须根除这些树干，才能植入美德。\OriginalTerm{斯多葛派}{Stoics}认为这些情欲必须被切除；\OriginalTerm{逍遥派}{Peripatetics}认为必须加以约束。两者都不正确，因为它们无法完全去除——它们是由本性植入的，并且具有确定而巨大的影响力；也无法削弱，因为如果它们是恶的，即使加以约束和适度使用，我们也应当没有它们；如果它们是善的，我们就应当完全使用它们。但我们说，它们既不应被去除，也不应被减弱。因为它们本身并非恶的，天主合理地将其植入我们内；而是因为它们在本质上是善的——它们被赐予我们是为了保护生命——却因滥用而变为恶。正如勇敢，如果你为保卫祖国而战，便是善；如果反对祖国，便是恶。同样，情欲，如果你用于善的目的，就是德行；用于恶，则被称为恶习。因此，愤怒是天主赐予的，用于抑制过犯，即控制属下的纪律，使恐惧抑制放纵、约束鲁莽。但不知其限度的人，会对同等的人甚至对上级发怒。由此导致残酷行为，由此导致杀戮，由此导致战争。贪财也被赐予，使我们渴望和寻求生活的必需品。但不知其边界的人，会贪得无厌地聚敛财富。由此产生毒害、欺诈、伪造遗嘱、各种欺骗。此外，情欲之欲是植入并内在于我们之中的，为了生育子女；但那些在心中不设定界限的人，只将其用于享乐。由此产生非法恋情、奸淫、放荡、各种败坏。因此，这些情欲必须保持在它们的界限内，并导向正确的方向；这样，即使它们强烈，也不会招致责备。\par

% structure: p0160 source=h2 target=section reason=relative-heading-rank; base=h2

\section{第六十二章 论节制感官的快乐}

当我们遭受伤害时，必须抑制愤怒，以便压制因争斗而迫近的邪恶，并持守两大美德：无害与忍耐。当我们拥有足够之物时，应克制贪欲。因为疯狂地劳作聚敛那些必将由抢劫、偷窃、抄家或死亡而转归他人的财物，是何等疯狂！情欲不应逾越婚姻之床，而应服务于生育子女。因为过度追求快乐既带来危险，也产生羞耻，尤其应当避免的是导致永死。没有什么比不洁的心灵和污秽的灵魂更令天主憎恶。也不要有人认为，只需戒绝这种快乐，即\OriginalTerm{从女性身体交合中获得的快乐}{quae capitur ex foeminei corporis copulatione}，还应戒绝由其他感官产生的快乐，因为它们本身也是败坏的，同一德行的一部分就是鄙弃它们。眼睛的快乐来自物体的美，耳朵的快乐来自和谐悦耳的声音，鼻子的快乐来自怡人的气味，味觉的快乐来自甜美的食物——所有这些，德行都应当强烈抵制，以免灵魂被这些诱惑所困，从天上的事物堕向地下的，从永恒堕向暂时，从永生堕向永罚。在味觉和嗅觉的快乐中，有这样的危险：它们能引我们走向奢侈。因为沉溺于这些的人，要么没有财产，要么即使有，也会挥霍一空，随后过上可憎的生活。但被听觉所诱的人（且不说歌曲，它们常常以某些精心编织的言辞、和谐的诗篇或巧妙的辩论，如此迷住内在感官，甚至以疯狂扰乱心灵的安定状态）容易转向不敬的崇拜。由此，那些能言善辩或偏好阅读优美文辞的人，便不易相信圣经，因为后者显得粗朴；他们不寻求真实之物，只寻求悦耳之物；甚至，在他们看来，最真实的就是那些悦耳的东西。于是他们拒绝真理，被话语的甜美所俘获。但涉及视觉的快乐是多样的。因为来自珍贵物品之美的快乐激起贪婪，这应当远离智者和义人；而来自女性容貌的快乐则将人引向另一种快乐，我们已在上文谈到。\par

% structure: p0162 source=h2 target=section reason=relative-heading-rank; base=h2

\section{第六十三章 论表演最能败坏心智}

现在需要讨论公共表演。由于这些表演对心灵的败坏有更强烈的影响，智者应当避免它们，并彻底防范，因为据说它们是为庆祝众神的荣耀而设立的。表演的展览是\Person{萨图尔努斯}的节庆。舞台属于\Person{利伯尔父神}；而车赛则被认为献给\Person{尼普顿}：因此，参与这些表演的人似乎已经离开了对天主的敬拜，转向了亵渎的仪式。但我更愿意谈论事情本身而非其起源。有什么比杀人更可怕、更可憎的呢？因此，我们的生命受到最严厉的法律保护；因此，战争是可憎的。然而，习俗却让人得以在无战争、无法律的情况下犯下杀人罪；而这对人来说是一种乐趣，因为他惩罚了罪行。但如果出席杀人现场意味着有罪，并且观众与行凶者承担同样的罪责，那么在这些角斗士的杀戮中，观众所沾染的血并不少于流血者；他既希望流血，又支持凶手并为其索要奖赏，就不能被视为未杀过人，也不能免于流血的罪责。至于舞台——难道它更神圣吗？舞台上，喜剧谈论荒淫与恋情，悲剧谈论乱伦与弑亲。演员的不雅姿态也模仿放荡女子，以舞蹈表达情欲。因为哑剧是败坏的学校，在那里，可耻之事通过象征性表现演出来，以便真事可以无所羞愧地发生。青年观看这些表演，他们处于危险的年龄，本应受到约束和管教，却通过这些表演被训练走向恶习与罪恶。至于竞技场，确实被认为较为无辜，但那里有更大的疯狂，因为观众的心灵如此疯狂，以至于他们不仅口出咒骂，还常常陷入争斗、打斗与冲突。因此，所有表演都应避免，以便我们能够保持心灵的安宁。我们必须弃绝有害的享乐，以免被致死的甜蜜所迷惑，陷入死亡的陷阱。\par

% structure: p0164 source=h2 target=section reason=relative-heading-rank; base=h2

\section{第六十四章——须制服情欲，戒绝禁物。}

惟愿美德取悦我们，其赏报是永生，当它战胜了享乐之时。但当情欲被征服、享乐被制服之后，对于追随天主与真理的人而言，压制其他事物就变得容易了：他永远不会咒骂，因为他盼望从天主那里得到祝福；他不会发假誓，以免嘲笑天主；他甚至不会起誓，以免因需要或习惯而陷入假誓。他不会说任何欺骗的话，不会说任何虚伪的话；他不会拒绝他已承诺的事，也不会承诺他无法完成的事；他不会嫉妒任何人，因为他对自己和自己的所有感到满足；他也不会夺走或希望他人遭受不幸——或许天主的恩惠更丰沛地赐予了那人。他不会偷窃，也不会贪图任何属于他人的东西。他不会将钱用于放债取利，因为那是以他人的不幸牟利；然而，如果某人因急需而必须借贷，他也不会拒绝出借。他对待儿子和奴仆不可苛刻：他应当记得自己也有一个父亲和主人。他对待他们的方式，就是他希望别人对待他的方式。他不会接受比自己资源更少的人所给的过度礼物；因为富人的产业不应通过穷人的损失而增加。\par

有一条古老的诫命：不可杀人。这条诫命不应被理解为仅仅禁止杀人——杀人甚至受到公共法律的惩罚。但藉着这条命令，我们也不得以言语致人于死地，也不得杀害或遗弃婴儿，也不得以自愿死亡来自我定罪。同样，我们被命令不可奸淫；但藉着这条诫命，我们不仅被禁止玷污他人的婚姻——这甚至受到万国公法的谴责——而且还要戒绝与卖淫者交往。因为天主的法律高于一切法律；它甚至禁止那些被视为合法的事，好使正义得以圆满。同一法律的一部分是不作假见证，而这条本身也有更广泛的含义。因为如果假见证以谎言伤害被控告者，并欺骗听者，那么我们就决不可说假话，因为谎言总是欺骗或伤害人。因此，即使不造成伤害，却口出虚言的人，也不是义人。他也不可奉承，因为奉承有害且欺诈；而他必处处持守真理。尽管这在当下可能令人不快，然而，当其益处与效用显现时，它不会像诗人所说的那样招致怨恨，而是带来感激。\par

% structure: p0167 source=h2 target=section reason=relative-heading-rank; base=h2

\section{第六十五章——关于命令之事与怜悯的训诫。}

我已经论述了那些被禁止的事；现在将简要说明哪些事是被命令的。与无害紧密相连的是怜悯。因为前者不施加伤害，后者则施行善工；前者开始正义，后者成全它。既然人的本性比其他动物更为软弱，而天主为后者提供了施暴的手段和抵御的防御，祂就赐予我们怜悯的情感，好使我们能将生命的全部保护置于相互帮助之中。因为我们是由同一天主所造，出于同一人，并由血缘之律相连，我们因此应当爱每一个人；所以我们不仅必须避免施加伤害，甚至当伤害加于我们时也不可报复，好使在我们内有无害的圆满。为此，天主命令我们甚至为仇敌时常祈祷。因此，我们应当成为一种适于陪伴与社会的动物，以便通过给予和接受帮助而相互保护。因为我们的脆弱易遭受许多意外和困难。要期望你所见发生在他人身上的事也可能临于你。这样，你若怀有那处于困苦中恳求你帮助者的心态，你终将被激发去施助。若有人缺乏食物，让我们施与；若有无衣者遇见我们，让我们给他衣服；若有人受强者欺凌，让我们解救他。让我们的房屋向陌生人或需要庇护者敞开。让我们的保护不缺少受监护者，让我们的护卫不缺少无助者。赎救俘虏是怜悯的伟大工作，探访和安慰贫穷的病患也是如此。若无助者或陌生人死去，我们不应任其不葬而暴露。这些是怜悯的工作和职责；若有人承担这些，他将向天主奉献真实而可悦纳的祭献。这祭品更适合奉献给天主，祂不因羊血而平息，却因人虔敬而喜悦；天主因是正义的，就以祂自己的法律和条件追随这人。祂向所见的怜悯者施怜悯；向所见的对恳求者苛刻者则毫不留情。因此，为能完成这一切悦乐天主的事，必须轻视钱财，并将其转移到天上的宝库里，那里贼不能挖透，锈不能腐蚀，暴君不能夺走，而是在天主的守护下为我们保存，直至我们永久的财富。\par

% structure: p0169 source=h2 target=section reason=relative-heading-rank; base=h2

\section{第66章——论宗教中的信德与刚毅。}

信德也是正义的一大组成部分；我们这些拥有信德之名者，尤其应在宗教中持守信德，因为天主先于且优于人。若为朋友、父母和子女，即为人而死是光荣的事，且行此事者获得长久的纪念与赞美，那么为天主而死，祂能以永生回报短暂的死亡，岂不更为光荣？因此，当此种必要临到，即我们被迫离弃天主，转向外邦人的礼仪时，没有恐惧、没有恐怖应使我们偏离所领受的信德。让天主在我们眼前，在我们心中，藉祂内在的帮助我们可克服肉体的痛苦和施加于身体的酷刑。那时，让我们只思念永生的赏报。如此，即使我们的肢体被撕裂或被焚烧，我们也能轻易承受暴虐残酷所能构想的一切。最后，让我们努力接受死亡本身，不是出于不情愿或胆怯，而是情愿而无畏地，如同那些知道在天主台前我们将获得何等光荣、已战胜世界并来到所应许给我们的事物面前的人，知道我们将以何等美善和莫大的福乐来补偿这些短暂的痛苦与今生所受的伤害。但若这光荣的机会缺乏，信德即使在平安中也必有赏报。\par

因此，让信德在生命的一切职责中、在婚姻中受到遵守。因为仅远离他人的床榻或妓院并不足够。凡有妻子者，勿再求他物，惟独满足于她，守护婚床的奧秘，使之纯洁无玷。因为在天主眼中，那卸下婚姻之轭，纵情于异性的欢愉，无论是自由女还是女奴，都是同等的奸淫与不洁。正如妇人受贞洁之链所束缚，不得渴慕其他男子，同样，丈夫也当受同一法律约束，因为天主将丈夫与妻子结合于一个身体。为此，祂命令除非妻子被判定犯奸，否则不得休妻，且婚姻契约的纽带永不解除，除非不忠已将其打破。为成全贞洁，还须加上：不仅行为上不犯罪，连意念中也不犯罪。因为显然，心意因欲望而被玷污，即使未实现；因此义人既不应行不义之事，也不应愿行不义之事。所以良心必须被洁净；因为不能被欺骗的天主洞察良心。胸膛必须清除一切污点，使之成为天主的宫殿，这宫殿不是由金银的光辉照亮，而是由信德与纯洁的光辉照亮。\par

% structure: p0172 source=h2 target=section reason=relative-heading-rank; base=h2

\section{第67章——论忏悔、灵魂不朽与天主眷顾。}

但这一切对人来说的确艰难，人性的软弱使任何人都不可能毫无瑕疵。因此，最后的补救办法便是投奔补赎；补赎在诸德之中绝非等闲之辈，因为它乃是自我的纠正；当我们偶然在言行上失足时，我们应立即回心转意，告明自己冒犯了天主，并恳求祂的宽恕；根据祂的仁慈，除非那些执迷不悟的人，祂必不拒绝赐予。补赎的帮助何其大，补赎的慰藉何其大！它是创伤与过犯的医治，是希望，是安全的港湾；谁若取消它，便是自绝救恩之路，因为没有人能如此正义，以致补赎对他永无必要。但我们，即使我们并无过犯，仍应向天主告明，为我们的缺失恳求宽恕，并在患难中也要感恩。让我们常向我们的主献上这份顺服。因为谦逊在天主眼中是可亲可爱的；既然祂更接纳那告明己罪的罪人，胜过傲慢的义人，那么祂岂不更接纳那告明的义人，并按其谦逊的程度在祂的天国中高举他？这些便是事奉天主者应当持守的；这便是祭品，这便是蒙悦纳的祭献；这便是真正的敬礼：人将他自己心灵的凭据献在天主的祭台上。至尊者喜悦这样的敬礼者，视之如子，并赐予他相称的永生赏报。关于永生，我现在必须谈论，并驳斥那些以为灵魂与身体一同消灭之人的主张。因为他们既不认识天主，又不能理解世界的奥秘，连人性和灵魂的本性也未参透。他们连要点都不掌握，又怎能看见后果？因此，在否认天主眷顾存在时，他们实际上是公然否认了天主——万有的泉源与根源。这导致他们要么主张现存之物永远存在，要么认为它们是自然产生或由微小种子的偶然结合而成。\par

不能说那存在且可见之物是永恒存在的，因为没有某种开端，它就不能凭自身存在。但没有任何事物能自然产生，因为没有不生自造者的本性。然而，最初如何能有种子呢？因为种子来自物体，而物体又来自种子。因此，没有无来源的种子。于是乎，当他们设想世界并非由天主眷顾所造时，他们甚至认为连人也不是按计划所造的。但若造人时未运用任何计划，那么灵魂便不能是不朽的。另一方面，另一些人认为只有一位天主，世界由祂所造，且为人而造，并且灵魂是不朽的。尽管他们持有真理性的观点，但他们仍未能洞察这一神圣工程与计划的原因、理由或结局，从而无法完成真理的全部奥秘并将其包含在某个界限之内。但他们因未把握完整的真理而未能做到的事，必须由我们——因天主的启示而知晓它的人——来完成。\par

% structure: p0175 source=h2 target=section reason=relative-heading-rank; base=h2

\section{第六十八章——论世界、人与天主眷顾。}

因此，让我们思考造就如此伟大工程的计划为何。天主创造了世界，正如\Person{柏拉图}所认为的，但他并未说明为何创造。他说，因为天主是善的，不嫉妒任何人，所以祂创造了善的事物。然而，我们看到自然体系中既有善亦有恶。或许会有某个乖戾的人站出来——如同那位无神论者\Person{塞奥多洛}一样——反驳\Person{柏拉图}说：不然，因为天主是恶的，所以祂创造了恶的事物。他又将如何驳斥呢？如果天主创造了善的事物，那么如此巨大的恶又从何而来？它们甚至常常胜过善的事物。对此，他说，恶包含在质料之中。但若曾有恶，则亦有善；因此，要么天主什么也没创造，要么若祂只创造了善的事物，那么那未被创造的恶事物就比那有开端的善事物更为永恒。因此，那些一度开始的事物将有终结，而那些永恒存在的事物将持久不变。因此，恶是可取的。但若恶并非可取，则它们的确不能更为永恒。因此，它们要么永恒存在，而天主无所作为；要么两者皆出于同一本源。因为天主创造万物，较之祂什么也没创造，更符合理性。因此，依据\Person{柏拉图}的观点，同一天主既是善的（因祂创造善事物），又是恶的（因祂创造恶事物）。如果这不可能，那么显然世界被天主创造的原因并非因为祂是善的。因为祂包含了善恶一切事物；祂并无为任何事物本身而创造，而是为了别的事物。建造房屋不仅是为了有房屋，而是为了接纳并庇护居住者。同样，造船不仅是为了看起来是一艘船，而是为了让人能在其中航行。器皿的制造也不仅是为了存在，而是为了容纳生活所需的物品。因此，天主创造世界也必有其用途。斯多葛派说它为人而造：此言不虚。因为人享用世界所包含的一切善物。但他们并未解释人本身为何被造，以及万物的创造者——天主眷顾——在人身上有何益处。\par

因此，依据\Person{柏拉图}的观点，同一天主既是善的（因祂创造善事物），又是恶的（因祂创造恶事物）。如果这不可能，那么显然世界被天主创造的原因并非因为祂是善的。因为祂包含了善恶一切事物；祂并无为任何事物本身而创造，而是为了别的事物。建造房屋不仅是为了有房屋，而是为了接纳并庇护居住者。同样，造船不仅是为了看起来是一艘船，而是为了让人能在其中航行。器皿的制造也不仅是为了存在，而是为了容纳生活所需的物品。因此，天主创造世界也必有其用途。斯多葛派说它为人而造：此言不虚。因为人享用世界所包含的一切善物。但他们并未解释人本身为何被造，以及万物的创造者——天主眷顾——在人身上有何益处。\par

\Person{柏拉图}也断言灵魂是不死的，但至于为什么、以何种方式、在何时、藉由谁的工具得以达至不死，或者那伟大奥秘的本性是什么，为何那些将要成为不死的人先前生于有死之身，然后，当他们完成了今世生命的历程、脱去脆弱身体的遮盖后，被转移到那永恒的幸福中——这一切他都无所领悟。最后，他没有解释天主的审判，也没有区分义人与不义之人，反而设想那些自陷于罪恶的灵魂至此受罚，即它们要在低级动物中重生，从而为罪行补赎，直到它们再次返回人形，并且这事永无止境，没有尽头。依我之见，他引入了某种类似梦境的戏谑，其中既没有天主的计划，也没有治理，更没有旨意。\par

% structure: p0179 source=h2 target=section reason=relative-heading-rank; base=h2

\section{第六十九章——世界是为了人而造的，人是为了天主而造的。}

我现在要说明那即使言说真理者也无法连贯起来、将原因与理由纳入同一视野的关键所在。世界由天主所造，为使人得以诞生；反之，人得以诞生，为要承认天主为父，在祂内有智慧；他们承认祂，为要崇拜祂，在祂内有正义；他们崇拜祂，为要领受不死的赏报；他们领受不死，为要永远事奉天主。你看见首者如何与中者紧密相连，中者如何与终者紧密相连吗？让我们分别审视它们，看它们是否彼此一致。天主为人之故创造了世界。凡看不见这一点的人，与畜类相差无几。除人以外，谁仰望天空？谁赞叹太阳、星辰以及天主的一切化工？谁居住在大地上？谁从大地收取果实？谁拥有鱼类、飞鸟、四足动物之权？岂非人么？因此，天主造了万物是为了人，因为万物都成为人的用处。\par

哲学家们看到了这一点，却没有看到进一步的结论：祂造人本身也是为了祂自己。因为既然祂为人造了如此伟大的工程，赐予人如此大的尊荣和权能，使人治理世界，那么人就应当承认如此宏恩的创造主——祂为人之故创造了世界本身——并应献上祂应得的崇拜和尊荣，这才是合宜、虔诚且必要的。\Person{柏拉图}在此失误了；在此他失去了最初把握的真理，因为他闭口不提他所承认的万有创造者与父的崇拜，也没有明白人因虔诚之纽带而连结于天主，宗教本身也由此得名，而灵魂之所以成为不死，正是为此。然而，他确实感知到灵魂是永恒的，但他没有通过有序的梯级下降到那一见解。因为撤去了中间的论证，他仿佛是陡然坠入真理，而非凭借理性偶然发现；他也没有更进一步，因为他是偶然而非凭理性找到真理。因此，应当崇拜天主，使人通过宗教（也就是正义）从天主那里领受不死，虔诚心灵没有其他赏报；如果这赏报是看不见的，则不能由看不见的天主赐予看得见的赏报。\par

% structure: p0182 source=h2 target=section reason=relative-heading-rank; base=h2

\section{第七十章——灵魂的不死得到证实。}

事实上，从许多论据可以收集到灵魂是永恒的。\Person{柏拉图}说，那永远自我运动、没有运动起始的，也没有终结；而人的灵魂总是自我运动，因为它能灵活思考、精巧发现、易于感知、适合学习，并且因为它记忆过去、理解现在、预见未来，通晓许多学科和技艺，所以它是不死的，因为它不包含任何与地上重量的沾染相混合之物。此外，灵魂的永恒性可以从德行和快乐上理解。快乐为一切动物共有，德行只属于人；前者是邪恶的，后者是高尚的；前者合乎本性，后者违反本性——除非灵魂是不死的。因为为信仰和正义辩护时，德行既不畏惧匮乏，也不因流放而惊慌，不惧怕监禁，不退缩于痛苦，也不拒绝死亡；既然这些事情违反本性，那么德行若不是愚蠢——因其阻碍利益、损害生命——那么灵魂就是不死的，它轻视今世之物，因为还有其他优于这些的事物，在身体解体后获得。但最大的不死证明是人独有对天主的认识。在哑巴动物中没有宗教观念，因为它们属土且俯向大地。人直立，仰望天空，正是为此，好寻求天主。因此，渴慕不死者不能不成为不死的。那与天主在面容和心智上相连的，不能败坏。最后，唯独人使用天上的元素，即火。因为如果光是通过火，生命是通过光，那么显然，使用火的人不是有死的，因为这火与那没有它便无光无生命者紧密相连、亲密相关。\par

但为何我们从论证推断灵魂是永生的，当我们有神圣的见证时？因为圣经和先知们的声音教导这一点。如果有人认为这不够，让他读\OriginalTerm{西比拉}{Sibyls}的诗篇，也让他考量\OriginalTerm{米利都的阿波罗}{Milesian Apollo}的答复，好让他明白\Person{德谟克利特}、\Person{伊壁鸠鲁}和\Person{狄凯阿科斯}是胡说，他们独独在所有人中否认那明显的事。证明灵魂的不朽之后，剩下的就是教导由谁、给谁、以何种方式、在何时赐予。既然固定且神定的时间开始被充满，万物的毁灭和终结必然发生，为的是世界能被天主更新。但那时刻近了，正如从年数和先知们预言的征兆所能收集到的。既然关于世界末日和时代终结所说的事情不可胜数，那些所说的事情本身必须不加修饰地陈述，因为提出证据将是无边无际的任务。如果有人想要它们，或不完全信任我们，让他接近天上经卷的神龛，并通过它们的可信度更充分地受教诲，让他明白哲学家们错了，他们认为这个世界是永恒的，或者从它被准备之时将有无数千年。因为六千年还未完成，当这个数目被凑满时，那时最终一切邪恶将被除去，为的是正义独自治。而这事将如何发生，我将用几句话解释。\par

% structure: p0185 source=h2 target=section reason=relative-heading-rank; base=h2

\section{第七十一章——论末后的时期。}

这些事由先知们说出，但作为观者，是即将要发生的。当最后的终结开始临近世界时，邪恶将增多；各种恶习和欺诈将变得频繁；正义将消失；信德、平安、仁慈、谦逊、真理将不复存在；暴力和胆大妄为将盛行；没有人会拥有任何东西，除非是用手获得，用手保卫。如果有任何好人，他们将被视为掠物和笑柄。没有人会对父母显孝爱，没有人会怜悯婴儿或老人；贪婪和情欲将败坏一切。将有屠杀和流血。将有战争，不仅是对外和邻邦之间的战争，还有内战。邦国之间将互相交战，每个性别和年龄都将拿起武器。政府的尊严将不被保持，军事纪律也不保持；但以强盗的方式，将有劫掠和蹂躏。王权将倍增，十个人将占领、分配和吞吃世界。将兴起另一个远更强大和邪恶的人，他毁灭了三个之后，将获得\Place{亚细亚}，并将其他列国征服制服在自己权力之下，将骚扰全地。他将制定新法律，废除旧法律；他将使国家成为自己的，并改变政府的名称和所在地。\par

那时将有一个可怕可憎的时期，没有人会选择活着。总之，事情的状态将是：哀号追随生者，庆贺追随死者。城市和乡镇将被摧毁，时而用火与剑，时而用反复的地震；忽而用水淹，忽而用瘟疫和饥荒。大地将不产任何东西，因严寒或酷热而荒芜。所有的水将部分变为血，部分因苦味而败坏，以致没有水可用于食物，或有益于饮用。在这些灾祸之上，还将添加来自天上的异象，使人们恐惧万状。彗星将频繁出现。太阳将被永恒的苍白遮蔽。月亮将被血染污，也不会弥补其被剥夺的光芒的损失。所有的星辰将坠落，季节也不会保持其规律，冬夏混淆。然后年、月和日都将缩短。\Person{特里斯墨吉斯忒斯}已宣告这是世界的老年和衰败。当这来临时，必须知道那时刻近了，天主将返回以改变世界。但在这些灾祸中，将兴起一个不虔敬的君王，不仅敌视人类，也敌视天主。他将践踏、折磨、骚扰并处死那些被先前暴君饶恕的人。那时将有不断的眼泪、永久的哀号和悲叹，以及对天主无用的祈祷；没有恐惧的休息，没有缓解的睡眠。白日总是增加灾难，黑夜则增加惊恐。如此，世界将减少到几乎荒无人烟，当然只剩下少数人。那时，那不虔敬的人也迫害义人和奉献给天主的人，并下令他自己应被当作神崇拜。因为他说他是基督，尽管他是祂的敌手。为使人相信，他将获得行奇迹的权能，以致火从天上降下，太阳离开其轨道，他设立的偶像能说话。通过这些异象，他将引诱许多人崇拜他，并在他们的手或额上接受他的记号。凡不崇拜他、不接受他记号的人，将死于残酷的折磨。如此，他将毁灭大约三分之二，三分之一将逃入荒凉的旷野。但他，疯狂且以不可平息的愤怒发怒，将率领军队围攻义人逃往的山。当他们看到自己被围困时，将以大声呼求天主的援助，天主将垂听他们，并派遣一位拯救者给他们。\par

% structure: p0188 source=h2 target=section reason=relative-heading-rank; base=h2

\section{第七十二章——论基督从天降下进行公审判，以及千年王国。}

那时，天将在暴风中开启，基督将带着大能降临，有火一样的亮光和无数天使队伍行在祂前面。那所有恶人的群众将被消灭，血流成河，而领袖本人将逃脱，并多次重整军队，第四次投入战斗；在此战斗中，他连同所有其他暴君被擒获，被交付火焚。但恶魔的王子本身，邪恶的创始者和策划者，将被火链捆绑，囚禁起来，好使世界得享和平，大地受多年蹂躏后得以安息。因此，和平建立，一切邪恶被压制后，那位正义的君王和得胜者要对地上活人和死人进行大审判，将使所有民族服从于活着的义人，并将义人死者复活得永生，祂自己将与他们一同在地上掌权，并建造圣城，这义人的王国将持续一千年。在那段时间，星辰将更加明亮，太阳的光辉将增强，月亮不再有亏缺。然后，从天主而来的祝福之雨将在早晚降下，大地将无需人的劳作而结出所有果实。蜂蜜将从岩石滴下，牛奶和酒泉将丰盈。野兽将放下凶猛变得温顺，狼在羊群中游荡而不伤害，牛犊与狮子一同吃草，鸽子与鹰隼结合，蛇不再有毒；没有动物靠流血存活。因为天主将为所有动物提供充足无害的食物。但当一千年期满，恶魔的王子被释放后，列国将反叛义人，不可胜数的人群将前来攻击圣人的城。那时，天主的最后审判将临到列国。因为祂将从根基震动大地，城市将被倾覆，祂将用硫磺与冰雹的火降在恶人身上，他们将着火，并互相残杀。但义人将暂时隐藏在地下，直到列国的毁灭完成，三天后他们出来，看见平原上布满尸首。然后将有地震，山岭崩裂，山谷深陷，死者的尸体将被堆积其中，其名称将称为\OriginalTerm{波吕安德里翁}{\GreekText{Πολυάνδριον}}。此后，天主将更新世界，将义人转变为天使的形体，使他们穿上不朽的衣袍，永远事奉天主；这就是天主的王国，没有终结。那时，恶人也复活，不是为生命，而是为惩罚；因为当第二次复活发生时，天主也将使他们复活，使他们被判定受永罚，交付永火，承受其罪行应得的惩罚。\par

% structure: p0190 source=h2 target=section reason=relative-heading-rank; base=h2

\section{第七十三章：得救的希望在于对天主的宗教和敬拜。}

因此，既然这一切都是真实而确定的，与先知预言宣告相符，既然\Person{特里斯墨吉斯托斯}、\Person{希斯塔斯普}和\Person{西彼拉}们早已预言了同样的事，那么所有生命和救恩的希望唯独在于对天主的宗教，这是无可置疑的。因此，除非一个人接受了\Person{基督}——天主已派遣并将为我们的救赎而派遣的祂，除非他通过\Person{基督}认识了至高天主，除非他遵守了祂的诫命和法律，否则他将落入我们所说的惩罚中。因此，必须轻视脆弱的事物，以便获得坚固的事物；必须藐视地上的事物，以便被天上的事物所尊荣；必须回避暂时的事物，以便达到永恒的事物。让每个人操练自己走向正义，塑造自己走向节制，预备自己走向竞赛，装备自己走向德行，以便万一敌人发动战争，他不会被任何力量、任何恐怖、任何折磨从正直和良善中驱离，不把自己交付给任何虚妄的虚构，而是在他的正直中承认唯一真天主，抛弃快乐——因快乐的诱惑，崇高的灵魂被压制到地上；坚持纯洁，尽可能帮助众人，藉善工为自己赚取不朽的财富，使他能以天主为审判者，为他的德行功绩获得信德的冠冕或不朽的赏报。\par

\SourceRecord{Translated by William Fletcher.；Ante-Nicene Fathers；Vol. 7.；Edited by Alexander Roberts, James Donaldson, and A. Cleveland Coxe.；Buffalo, NY: Christian Literature Publishing Co.,；1886.；<http://www.newadvent.org/fathers/0702.htm>.}
